\documentclass[11pt,a4paper,openany]{report}

\usepackage{fontspec}
\usepackage{mathtools}
\usepackage{unicode-math}
\usepackage[inner=27mm,outer=24mm,top=26mm,bottom=25mm,headheight=23pt]{geometry}
\usepackage{xcolor}
\usepackage{hyperref}
\usepackage{fancyhdr}
\usepackage{titlesec}
\usepackage{enumitem}
\usepackage{cite}
\usepackage{amsthm}
\usepackage{booktabs}
\usepackage{longtable}
\usepackage{array}
\usepackage{graphicx}

\setmainfont[
  Path=fonts/,
  BoldFont=NotoSerifJP-Bold.ttf,
  ItalicFont=NotoSerifJP-Regular.ttf,
  ItalicFeatures={FakeSlant=0.15},
  BoldItalicFont=NotoSerifJP-Bold.ttf,
  BoldItalicFeatures={FakeSlant=0.15}
]{NotoSerifJP-Regular.ttf}
\setsansfont[
  Path=fonts/,
  BoldFont=NotoSansJP-Bold.ttf,
  ItalicFont=NotoSansJP-Regular.ttf,
  ItalicFeatures={FakeSlant=0.15},
  BoldItalicFont=NotoSansJP-Bold.ttf,
  BoldItalicFeatures={FakeSlant=0.15}
]{NotoSansJP-Regular.ttf}
\setmonofont{Latin Modern Mono}
\setmathfont{Latin Modern Math}

\XeTeXlinebreaklocale "ja"
\XeTeXlinebreakskip=0pt plus 1pt
\emergencystretch=2em
\allowdisplaybreaks[2]
\numberwithin{equation}{chapter}
\setcounter{tocdepth}{2}

\definecolor{PaperNavy}{HTML}{173B57}
\definecolor{PaperBlue}{HTML}{275D78}
\definecolor{PaperRule}{HTML}{9AB4C3}
\definecolor{PaperTint}{HTML}{F2F7FA}

\hypersetup{
  unicode=true,
  colorlinks=true,
  linkcolor=PaperBlue,
  citecolor=PaperBlue,
  urlcolor=PaperBlue,
  pdftitle={有限正準位相担体の相関行列力学：共通試行測度による有限基底 Born 型標本化と Bell 型履歴統計},
  pdfauthor={Open research manuscript}
}

\titleformat{\chapter}[display]
  {\normalfont\sffamily\bfseries\color{PaperNavy}\centering}
  {\Large 第\thechapter 章}
  {0.8ex}
  {\Huge}
  [\vspace{1.1ex}\color{PaperRule}\titlerule]
\titleformat{name=\chapter,numberless}[display]
  {\normalfont\sffamily\bfseries\color{PaperNavy}\centering}
  {}{0pt}{\Huge}
  [\vspace{1.1ex}\color{PaperRule}\titlerule]
\titlespacing*{\chapter}{0pt}{-8pt}{28pt}

\titleformat{\section}
  {\normalfont\Large\sffamily\bfseries\color{PaperNavy}}
  {\thesection}{0.7em}{}
\titleformat{\subsection}
  {\normalfont\large\sffamily\bfseries\color{PaperBlue}}
  {\thesubsection}{0.7em}{}
\titlespacing*{\section}{0pt}{2.8ex plus 0.6ex minus 0.2ex}{1.2ex}
\titlespacing*{\subsection}{0pt}{2.3ex plus 0.5ex minus 0.2ex}{0.9ex}

\pagestyle{fancy}
\fancyhf{}
\fancyhead[L]{\small\color{PaperBlue}有限正準位相担体の相関行列力学}
\fancyhead[R]{\small\color{PaperBlue}\nouppercase{\leftmark}}
\fancyfoot[C]{\thepage}
\renewcommand{\headrulewidth}{0.4pt}
\renewcommand{\headrule}{\hbox to\headwidth{\color{PaperRule}\leaders\hrule height \headrulewidth\hfill}}
\renewcommand{\chaptermark}[1]{\markboth{\thechapter\quad #1}{}}

\setlength{\parindent}{1em}
\setlength{\parskip}{0.35em}
\linespread{1.08}
\clubpenalty=10000
\widowpenalty=10000
\displaywidowpenalty=10000
\setlist[itemize]{leftmargin=2.2em,itemsep=0.25em,topsep=0.45em}
\setlist[enumerate]{leftmargin=2.4em,itemsep=0.25em,topsep=0.45em}

\newenvironment{statusbox}{%
  \par\medskip\noindent\begin{center}\begin{minipage}{0.92\linewidth}%
  \color{PaperBlue}\small\sffamily\hrule height 0.6pt\vspace{0.7em}%
}{%
  \vspace{0.7em}\hrule height 0.6pt\end{minipage}\end{center}\medskip%
}

\newtheorem{theorem}{定理}[chapter]
\newtheorem{proposition}[theorem]{命題}
\newtheorem{lemma}[theorem]{補題}
\newtheorem{corollary}[theorem]{系}
\theoremstyle{definition}
\newtheorem{assumption}[theorem]{仮定}
\newtheorem{definition}[theorem]{定義}
\theoremstyle{remark}
\newtheorem{remark}[theorem]{注意}
\renewcommand{\proofname}{証明}

\newcommand{\tightlist}{\setlength{\itemsep}{0pt}\setlength{\parskip}{0pt}}
\renewcommand{\contentsname}{目次}
\renewcommand{\bibname}{参考文献}
\renewcommand{\figurename}{図}
\renewcommand{\tablename}{表}

\makeatletter
\renewcommand*\@pnumwidth{2.2em}
\renewcommand*\@tocrmarg{3.2em}
\renewcommand*\l@section{\@dottedtocline{1}{1.5em}{3.2em}}
\renewcommand*\l@subsection{\@dottedtocline{2}{4.7em}{3.8em}}
\makeatother

\begin{document}

\begin{titlepage}
  \centering
  \vspace*{29mm}
  {\sffamily\bfseries\color{PaperNavy}\fontsize{20}{29}\selectfont
    有限正準位相担体の相関行列力学：\\[2mm]
    共通試行測度による\\[2mm]
    有限基底 Born 型標本化と\\[2mm]
    Bell 型履歴統計\par}
  \vspace{9mm}
  {\Large\color{PaperBlue}相関行列、正準作用選択器、\\[1mm]
    設定条件付き完全履歴\par}
  \vspace{14mm}
  {\color{PaperRule}\rule{0.68\textwidth}{0.9pt}\par}
  \vfill
  {\large draft-21 改訂稿\par}
  \vspace{22mm}
\end{titlepage}

\pagenumbering{roman}
\tableofcontents
\clearpage
\pagenumbering{arabic}

\chapter*{概要}

\addcontentsline{toc}{chapter}{概要}

本論文は、有限自由度の古典 Hamiltonian 系を基礎とし、対象部分を外部との微小なエネルギー交換を伴う弱開放系として扱う。目的は、量子力学をミクロ構成の入力に置かず、量子力学に特徴的な力学と確率構造が有効理論として現れ得る範囲を明らかにすることである。

各試行で実在するものを、粒子、有限個の正準位相担体、源、浴、時計、測定器、比較器、記録装置、外部自由度に限定する。全自由度の位相空間を \(\Gamma_{\rm all}\)、自律 Hamiltonian を \(H_{\rm all}\)、流れを \(\Phi_t\) とする。時計の正方向横断で定めた試行開始面 \(\Sigma_0\) と、次の開始面までの周期写像 \(\mathcal R\) を導入する。本稿は、\(\Sigma_0\) 上に不変かつエルゴード的な母測度 \(\mu_*\) が存在し、ほとんど全ての開始点が有限時間で次の試行へ戻ることを中心仮定とする。エルゴード性により、同じ装置を繰り返した長期頻度を \(\mu_*\) の事象確率と同一視できる。ただし、具体的な全 Hamiltonian から目的の \(\mu_*\) を導くことは未完成である。

実正準対 \((Q_i,P_i)\) から、固定作用尺度 \(\mathcal J_0>0\) を用いて

\[
b_i
=
\frac{Q_i+iP_i}{\sqrt{2\mathcal J_0}}
\]

を定める。複素数は量子的公理ではなく、実2次元正準平面の表示である。調製条件 \(\mathcal P\) と測定プログラム \(M\) で条件付けた相関行列を

\[
C(t)
=
\mathbb E_{\mu_*}
\left[
b_t b_t^\dagger
\mid
\mathcal P,M
\right]
\]

とする。\(C\) は正半定値 Hermitian 行列であり、母測度の時刻別2次モーメントである。各試行で \(C\) という物質が存在するとは解釈しない。

第1の主結果は、有限2次 Hamiltonian

\[
H_{\rm ph}
=
b^\dagger h_L(t)b,
\qquad
h_L(t)=h_L(t)^\dagger
\]

から、各試行の振幅方程式と相関行列方程式

\[
i\mathcal J_0\dot b
=
h_Lb,
\qquad
i\mathcal J_0\dot C
=
\left[h_L,C\right]
\]

が厳密に従うことである。同じ条件付き集団の全試行が共通の \(h_L\) に従う必要がある。閉鎖線形発展では \(C\) の跡、固有値、階数、純度が保存される。従って、混合集団が観測中に自発的に階数1へ純化することはない。共通源が階数1相関 \(C=\Lambda\chi\chi^\dagger\) を準備した場合に限り、共通位相を除いて \(i\mathcal J_0\dot\chi=h_L\chi\) を得る。

第2の主結果は、正準作用選択器による有限基底 Born 型標本化である。測定基底を表すユニタリ正準混合を \(U\) とし、各試行の出力作用を

\[
I_k
=
\mathcal J_0
\left|
\left(Ub\right)_k
\right|^2,
\qquad
I_{\rm ph}
=
\sum_k I_k
\]

とする。位相担体、基底、調製条件の下で条件付き一様な正準角から \([0,I_{\rm ph})\) 上の選択変数を作り、長さ \(I_k\) の累積区間を1つ選ぶ。このとき一般に

\[
P(k\mid U,\mathcal P)
=
\mathbb E_{\mu_*}
\left[
\frac{I_k}{I_{\rm ph}}
\middle|
U,\mathcal P,M=\mathsf{basis}
\right]
\]

が厳密に成立する。さらに全位相作用が試行ごとに固定され、調製相関が \(U\) に依存しなければ、

\[
P(k\mid U,\mathcal P)
=
\frac{
\left(UC U^\dagger\right)_{kk}
}{
\operatorname{tr}C
}
\]

を得る。階数1は不要であり、高階数相関行列も扱える。\(U=I\) でモードが空間セルなら位置基底の重みになるが、これは位相担体モードの選択であり、粒子の局所検出または連続軌道を自動的には与えない。

第3の主結果は、同じ母測度から取り出す設定条件付き Bell 完全履歴統計である。設定記号を \(x,y\)、局所結果を \(A,B\in\{+1,-1\}\) とする。反対称左右交差相関と局所分析器から

\[
K_{AB}^{xy}
=
\frac{\mathcal K}{4}
\left[
1-AB\cos\Delta_{xy}
\right]
\]

という非負枝作用が厳密に得られる。局所結果は共通未来より前に記録されるため、未来の選択器を結果生成器とは解釈しない。共通未来の正準角は、既に記録された結果と枝区間が整合する完結履歴を定める。物理的母測度が整合履歴だけに支持を持ち、4結果の基準セクター密度が共通なら、

\[
P(A,B\mid x,y)
=
\frac14
\left[
1-AB\cos\Delta_{xy}
\right]
\]

が従う。理想対称条件では非信号周辺と CHSH 値 \(2\sqrt2\) を得る。設定生成器も母測度に含まれ、一般に \(\mu_*(d\Lambda\mid x,y)\neq\mu_*(d\Lambda)\) となる。Bell の前提違反は測定設定独立性にあり、Bell の定理を否定しない。未来の比較器が過去の局所記録を変更する機構も導入しない。

3つの統計は別々の基礎測度ではなく、\(\mu_*\) の2次モーメント、有限基底結果事象、設定条件付き完結履歴事象として統一される。ただし、位置側と Bell 側の取り出し方は同じではない。位置側の正準角は前向きな結果選択器である一方、Bell 側の正準角は二側履歴の整合変数である。

本文は6章から成る。第1章で統一試行系、母測度、主張の範囲、第2章で有限正準位相担体と相関行列、第3章で相関行列力学、空間グラフ、干渉と節、第4章で正準作用選択器と有限基底 Born 型標本化、第5章で Bell 余弦枝作用と設定条件付き履歴測度、第6章で誤差、反証条件、未解決問題を扱う。正準計算、母測度の条件付け、選択器定理、装置 Hamiltonian の候補は付録へ置く。

最大の未解決問題は4つある。第1に、局所古典結合から Schrödinger 型係数を設計値でなく必然的に得ること。第2に、累積レジスター、比較器、唯一結果形成、記録、逆計算を1つの滑らかな有限 Hamiltonian 周期へ統合すること。第3に、位相担体の有限基底選択を粒子の局所検出と連続運動へ接続すること。第4に、Bell 整合支持条件と共通基準セクター密度を持つ不変母測度を、事後選別なしの具体的な周期写像から導くことである。

\hypertarget{ux7d71ux4e00ux8a66ux884cux7cfbux6bcdux6e2cux5ea6ux4e3bux7d50ux679c}{%
\chapter{統一試行系、母測度、主結果}\label{ux7d71ux4e00ux8a66ux884cux7cfbux6bcdux6e2cux5ea6ux4e3bux7d50ux679c}}

\begin{statusbox}
位置づけ：相関行列、有限基底結果、Bell完全履歴統計を1つの不変母測度から定義する。不変性とエルゴード性は仮定する。正準作用選択器の確率式は条件付き一様性の下で厳密であり、Bell共同確率には二側履歴整合と共通基準セクター密度が必要である。
\end{statusbox}

\hypertarget{ux554fux984cux8a2dux5b9a}{%
\section{問題設定}\label{ux554fux984cux8a2dux5b9a}}

本論文の目的は、明示的な古典 Hamiltonian 系から、量子力学に特徴的な確率構造が有効理論として出現し得るかを検証することである。量子力学をミクロ Hamiltonian の入力には使わず、得られた式と測定統計の特徴を判定するときだけ参照する。

現行稿は、以前は別々に置いていた3つの統計構造を1つの試行空間へ統合する。

\begin{enumerate}
\def\labelenumi{\arabic{enumi}.}
\tightlist
\item
  相関行列は、母測度の時刻別2次モーメントである。
\item
  有限基底結果は、条件付き一様な正準角が選ぶ累積作用区間である。
\item
  Bell統計は、設定で条件付けた完結履歴事象の頻度である。
\end{enumerate}

統一とは、3つを同じ確率公式へ還元することではない。同じ母測度から異なる観測写像で取り出すことを意味する。位置側の選択器は結果形成より前に働けるが、Bell側の共通未来変数は、既に形成された局所結果を作り直せない。この時間順序の違いを保つ。

\hypertarget{ux62e1ux5927ux5168ux7cfbux3068ux30d7ux30edux30b0ux30e9ux30e0ux5909ux6570}{%
\section{拡大全系とプログラム変数}\label{ux62e1ux5927ux5168ux7cfbux3068ux30d7ux30edux30b0ux30e9ux30e0ux5909ux6570}}

全正準状態を概略

\[
Z
=
\left(
X,P_X,
b,
Y_{\rm src},
Y_{\rm set},
Y_{\rm det},
Y_{\rm cmp},
Y_{\rm mem},
Y_{\rm clk},
Y_{\rm ext}
\right)
\]

と書き、全位相空間を \(\Gamma_{\rm all}\) とする。\((X,P_X)\) は粒子、\(b\) は有限位相担体、残る変数は源、設定生成器、検出器、比較器、記録、時計、外部環境を表す。自律 Hamiltonian と流れを

\[
H_{\rm all}(Z),
\qquad
\Phi_t:
\Gamma_{\rm all}
\longrightarrow
\Gamma_{\rm all}
\]

とする。有限装置部分だけを見ると、記録情報とエネルギーが外部へ移り得る。外部自由度まで含む拡大全系では Hamiltonian 流として扱う。

測定プログラムを離散レジスター \(M\) で表し、相関観測、有限基底測定、Bell測定を区別する。離散値は、互いに分離した正準レジスター領域の略記であり、Hamiltonian が不連続であることを意味しない。単粒子測定基底を \(U\)、Bell設定を \(x,y\) とする。位相担体 \(b\) とBell右設定の記号を重ねない。

\hypertarget{ux8a66ux884cux958bux59cbux9762ux3068ux5468ux671fux5199ux50cf}{%
\section{試行開始面と周期写像}\label{ux8a66ux884cux958bux59cbux9762ux3068ux5468ux671fux5199ux50cf}}

時計角を \(\vartheta_{\rm clk}\) とする。試行開始面を

\[
\Sigma_0
=
\left\{
z:
\vartheta_{\rm clk}=0,
\dot\vartheta_{\rm clk}>0
\right\}
\cap
\Gamma_{\rm prep}
\]

と定める。\(\Gamma_{\rm prep}\) は、源、指針、比較器が所定の準備範囲にあり、外部記録と不要情報には多数のミクロ状態を許す準備マクロ領域である。全自由度を同じ1点へ戻す面ではない。

最初の再帰時刻と周期写像を

\[
\tau(z)
=
\inf
\left\{
t>T_{\min}:
\Phi_tz\in\Sigma_0
\right\},
\]

\[
\mathcal R(z)
=
\Phi_{\tau(z)}z
\]

とする。結果別に異なる再横断を重複して数えないため、\(T_{\min}>0\) は1試行の測定窓より長く取る。戻らない周期を黙って標本から外さず、

\[
\mu_*
\left(
\tau<\infty
\right)
=
1
\]

を要求する。

\hypertarget{ux4e0dux5909ux304bux3064ux30a8ux30ebux30b4ux30fcux30c9ux7684ux306aux6bcdux6e2cux5ea6}{%
\section{不変かつエルゴード的な母測度}\label{ux4e0dux5909ux304bux3064ux30a8ux30ebux30b4ux30fcux30c9ux7684ux306aux6bcdux6e2cux5ea6}}

\(\Sigma_0\) 上の確率測度 \(\mu_*\) に対し、

\[
\mathcal R_*\mu_*
=
\mu_*
\]

を仮定する。さらに本稿では、\(\mathcal R\) が \(\mu_*\) に関してエルゴード的であると仮定する。従って任意の可積分な試行観測量 \(f\) に対し、典型的な開始点で

\[
\frac1N
\sum_{n=0}^{N-1}
f\left(\mathcal R^nz\right)
\longrightarrow
\int_{\Sigma_0}
f
\,d\mu_*
\]

が成立する。事象指示関数を選べば、単一装置の長期頻度が \(\mu_*\) の事象確率へ収束する。

不変性とエルゴード性は現行模型の仮定であり、全 Hamiltonian から導いた結論ではない。エルゴード性は収束率を与えないため、有限試行数の自己相関と混合時間は別に評価する。また、目的の \(\mu_*\) の存在、一意性、物理的準備は未解決である。

\hypertarget{ux3064ux306eux7d71ux8a08ux306eux5b9aux7fa9}{%
\section{3つの統計の定義}\label{ux3064ux306eux7d71ux8a08ux306eux5b9aux7fa9}}

開始点 \(z\in\Sigma_0\) から時刻 \(t\) の位相担体を

\[
b_t(z)
=
b\left(\Phi_tz\right)
\]

とする。調製条件 \(\mathcal P\) とプログラム \(M\) で条件付けた相関行列は

\[
C_M(t)
=
\mathbb E_{\mu_*}
\left[
b_t b_t^\dagger
\mid
\mathcal P,M
\right]
\]

である。

有限基底測定で結果レジスターが \(k\) を記録する事象を \(E_k^U\) とする。結果頻度は

\[
P(k\mid U,\mathcal P)
=
\mu_*
\left(
E_k^U
\mid
U,\mathcal P,M=\mathsf{basis}
\right)
\]

である。

Bell設定生成器と局所記録を、それぞれ \(x(z),y(z)\) と \(A(z),B(z)\) とする。完結履歴事象を

\[
F_{AB}^{xy}
=
\left\{
z:
x(z)=x,
y(z)=y,
A(z)=A,
B(z)=B
\right\}
\]

と定めれば、

\[
P(A,B\mid x,y)
=
\frac{
\mu_*
\left(
F_{AB}^{xy}
\cap
\{M=\mathsf{Bell}\}
\right)
}{
\mu_*
\left(
\{x(z)=x,y(z)=y,M=\mathsf{Bell}\}
\right)
}
\]

である。分母が正の設定対だけを扱う。相関行列、有限基底頻度、Bell頻度の基礎測度は全て同じ \(\mu_*\) である。

\hypertarget{ux4e3bux5b9aux74061ux306eux7bc4ux56f2}{%
\section{主定理1の範囲}\label{ux4e3bux5b9aux74061ux306eux7bc4ux56f2}}

第2章と第3章の主定理1は、次をまとめた有限次元結果である。

\begin{enumerate}
\def\labelenumi{\arabic{enumi}.}
\tightlist
\item
  \(H_{\rm ph}=b^\dagger h_Lb\) から \(i\mathcal J_0\dot b=h_Lb\) が厳密に従う。
\item
  同じ条件付き集団の全試行が共通 \(h_L\) に従うなら、\(i\mathcal J_0\dot C_M=[h_L,C_M]\) が厳密に閉じる。
\item
  閉鎖線形発展では \(C_M\) の固有値、階数、純度が保存される。
\item
  共通源が階数1相関を準備すれば、その因子 \(\chi\) は共通位相を除いて Schrödinger 型方程式に従う。
\item
  近似階数1の欠陥は干渉節の残留強度を上から抑える。
\end{enumerate}

プログラム、設定、ミクロ状態によって \(h_L\) が試行ごとに変わる場合は、単純な交換子閉包は一般に成立しない。測定と再初期化をまたぐ全周期が同じ交換子方程式に従うとも主張しない。

\hypertarget{ux4e3bux5b9aux74062ux306eux7bc4ux56f2}{%
\section{主定理2の範囲}\label{ux4e3bux5b9aux74062ux306eux7bc4ux56f2}}

第4章の主定理2は、正準作用選択器による有限基底標本化である。各試行で測定前の正準混合 \(U\) を施し、

\[
I_k
=
\mathcal J_0
\left|
\left(Ub\right)_k
\right|^2,
\qquad
I_{\rm ph}
=
\sum_kI_k
\]

を作る。選択器角が \((b,U,\mathcal P)\) の下で条件付き一様なら、

\[
P(k\mid U,\mathcal P)
=
\mathbb E_{\mu_*}
\left[
\frac{I_k}{I_{\rm ph}}
\middle|
U,\mathcal P,M=\mathsf{basis}
\right]
\]

が厳密に成立する。\(I_{\rm ph}\) が試行ごとに固定され、調製相関が \(U\) に依存しなければ、

\[
P(k\mid U,\mathcal P)
=
\frac{
\left(UC U^\dagger\right)_{kk}
}{
\operatorname{tr}C
}
\]

を得る。階数1は不要である。全作用が変動する場合は比の平均と平均の比を区別し、共分散補正を残す。

本定理は、条件付き一様角、累積作用読出し、比較、唯一結果記録が働くことを仮定する。候補となる正準レジスターは付録Cで示すが、全てを連結した滑らかな完成 Hamiltonian は未構成である。

\hypertarget{ux4e3bux5b9aux74063ux306eux7bc4ux56f2}{%
\section{主定理3の範囲}\label{ux4e3bux5b9aux74063ux306eux7bc4ux56f2}}

第5章の主定理3は、同じ母測度上の設定条件付き Bell 完全履歴統計である。論理の順序は次である。

\begin{enumerate}
\def\labelenumi{\arabic{enumi}.}
\tightlist
\item
  反対称左右交差相関と局所分析器から余弦枝作用 \(K_{AB}^{xy}\) を得る。
\item
  左右の局所結果 \(A,B\) は、戻り信号が共通未来へ到達する前に記録される。
\item
  共通未来で枝作用と正準角から4つの区間を計算する。
\item
  未来角が、既に記録された \((A,B)\) の区間に入る完結履歴だけを物理的支持とする。
\item
  整合前の4結果セクターの基準密度が共通なら、余弦共同確率を得る。
\item
  対称条件から非信号周辺と CHSH 値を得る。
\item
  設定条件付き母測度により、Bell の測定設定独立性が成立しないことを特定する。
\end{enumerate}

未来選択器は結果生成器でなく履歴整合変数である。共通基準セクター密度と整合支持条件を具体的な周期写像から導くことは、主定理3の仮定であって結論ではない。観測条件付け \(P(A,B\mid x,y)\) と外部介入で設定だけを変更した分布も同一視しない。

\hypertarget{ux5c0eux51faux72b6ux614b}{%
\section{導出状態}\label{ux5c0eux51faux72b6ux614b}}

\begin{longtable}[]{@{}
  >{\raggedright\arraybackslash}p{(\columnwidth - 4\tabcolsep) * \real{0.3333}}
  >{\raggedright\arraybackslash}p{(\columnwidth - 4\tabcolsep) * \real{0.3333}}
  >{\raggedright\arraybackslash}p{(\columnwidth - 4\tabcolsep) * \real{0.3333}}@{}}
\toprule\noalign{}
\begin{minipage}[b]{\linewidth}\raggedright
主張
\end{minipage} & \begin{minipage}[b]{\linewidth}\raggedright
導出状態
\end{minipage} & \begin{minipage}[b]{\linewidth}\raggedright
主な制限
\end{minipage} \\
\midrule\noalign{}
\endhead
\bottomrule\noalign{}
\endlastfoot
実正準対から複素位相担体への表示 & 厳密結果 & 有限自由度 \\
相関行列の交換子発展 & 厳密結果 & 母測度の同じ条件付き集団で共通の2次 Hamiltonian \\
跡、固有値、階数、純度の保存 & 厳密結果 & 閉鎖線形発展 \\
共通源による階数1準備 & 厳密結果 & 理想正準準備回路 \\
階数1因子の Schrödinger 型発展 & 厳密結果・仮説依存 & 階数1準備と \(h_L\) の選択 \\
正準作用選択器の一般確率式 & 厳密結果・仮説依存 & 条件付き一様角と唯一結果形成 \\
固定作用下の有限基底 Born 型則 & 厳密結果・仮説依存 & 固定全作用と単粒子側の基底独立準備 \\
Bell余弦枝作用 & 厳密結果・仮説依存 & 反対称源と局所平面回転 \\
共通基準密度下の Bell履歴統計 & 厳密結果・仮説依存 & 整合支持条件と基準密度の物理的準備は未完成 \\
長期頻度と母測度確率の一致 & 厳密結果・仮説依存 & 周期写像の不変性とエルゴード性を仮定 \\
\end{longtable}

\hypertarget{ux672cux8ad6ux6587ux304cux4e3bux5f35ux3057ux306aux3044ux3053ux3068}{%
\section{本論文が主張しないこと}\label{ux672cux8ad6ux6587ux304cux4e3bux5f35ux3057ux306aux3044ux3053ux3068}}

本論文は次を主張しない。

\begin{enumerate}
\def\labelenumi{\arabic{enumi}.}
\tightlist
\item
  目的の母測度の存在、一意性、不変性、エルゴード性を具体的な全 Hamiltonian から導いたこと。
\item
  任意初期集団からの階数1純化。
\item
  古典振動子による Schrödinger 方程式表示そのものの新規性。
\item
  \(h_L\) の係数と局所結合をより原始的な古典模型から必然的に導いたこと。
\item
  非線形ミクロ Hamiltonian に対する無条件な相関行列閉包。
\item
  位相担体の有限基底選択が粒子の局所検出そのものであること。
\item
  粒子の連続軌道または全時刻での \(|\psi|^2\) 分布。
\item
  Wallstrom 問題の解決または連続空間での循環量子化。
\item
  連続スペクトルを含む一般 Born 則。
\item
  一般の複合状態または一般 Tsirelson 原理。
\item
  Bell共通基準密度と履歴整合支持を生成する完成周期。
\item
  Bell の定理の否定。
\item
  後の比較器が過去の記録を変更する逆因果制御。
\item
  二側境界模型と前周期記憶による共通原因模型の同値性。
\end{enumerate}

配置拡散・Nelson経路、旧2成分誘導場、旧位置作用殻による標本化、前周期記憶模型は現行主線に採用しない。再利用価値のある結果と不採用理由は研究メモまたは版履歴で管理する。

\hypertarget{ux6709ux9650ux6b63ux6e96ux4f4dux76f8ux62c5ux4f53ux3068ux968eux65701ux76f8ux95a2ux306eux6e96ux5099}{%
\chapter{有限正準位相担体と階数1相関の準備}\label{ux6709ux9650ux6b63ux6e96ux4f4dux76f8ux62c5ux4f53ux3068ux968eux65701ux76f8ux95a2ux306eux6e96ux5099}}

\begin{statusbox}
位置づけ：実正準対、複素振幅、保存全位相作用、相関行列、階数1の同値条件、共通源準備は有限次元で厳密である。一般集団の純化、源の反復可能な再初期化、非線形閉包は未完成である。
\end{statusbox}

\hypertarget{ux5b9fux6b63ux6e96ux5bfeux3068ux8907ux7d20ux632fux5e45}{%
\section{実正準対と複素振幅}\label{ux5b9fux6b63ux6e96ux5bfeux3068ux8907ux7d20ux632fux5e45}}

有限モード \(i=1,\ldots,L\) に実正準対

\[
\left(Q_i,P_i\right),
\qquad
\left\{Q_i,P_j\right\}
=
\delta_{ij}
\]

を置く。固定作用尺度 \(\mathcal J_0>0\) を用いて

\[
b_i
=
\frac{Q_i+iP_i}{\sqrt{2\mathcal J_0}}
\]

と定める。Poisson 括弧は

\[
\left\{
b_i,b_j^*
\right\}
=
-
\frac{i}{\mathcal J_0}
\delta_{ij}
\]

である。\(b_i\) は実2次元正準平面の表示であり、量子的な生成消滅演算子ではない。

各モードの作用と全位相作用を

\[
I_i
=
\frac12
\left(
Q_i^2+P_i^2
\right)
=
\mathcal J_0
\left|b_i\right|^2,
\]

\[
I_{\rm ph}
=
\sum_iI_i
=
\mathcal J_0b^\dagger b
\]

とする。\(I_{\rm ph}\) は状態に依存する保存量、\(\mathcal J_0\) は座標尺度である。

\hypertarget{ux6709ux96502ux6b21-hamiltonian}{%
\section{有限2次 Hamiltonian}\label{ux6709ux96502ux6b21-hamiltonian}}

有限 Hermitian 行列 \(h(t)=h(t)^\dagger\) に対して

\[
H_{\rm ph}(t)
=
b^\dagger h(t)b
\]

と置く。この量は実数であり、\(Q_i,P_i\) だけからなる通常の2次 Hamiltonian である。

\begin{theorem}[有限正準位相担体の発展]
$H_{\rm ph}=b^\dagger hb$ の Hamilton 方程式は

$$
i\mathcal J_0\dot b
=
hb
$$

である。$h$ が Hermitian なら $I_{\rm ph}$ は保存される。
\end{theorem}

\begin{proof}
複素 Poisson 括弧から

$$
\dot b_i
=
\left\{b_i,H_{\rm ph}\right\}
=
-
\frac{i}{\mathcal J_0}
\frac{\partial H_{\rm ph}}{\partial b_i^*}
=
-
\frac{i}{\mathcal J_0}
\left(hb\right)_i
$$

を得る。また

$$
\frac{d}{dt}
\left(b^\dagger b\right)
=
\frac{i}{\mathcal J_0}
b^\dagger hb
-
\frac{i}{\mathcal J_0}
b^\dagger hb
=
0
$$

である。
\end{proof}

共通回転 \(b\mapsto e^{i\beta}b\) は全位相作用 \(I_{\rm ph}\) が生成する。観測可能な混合と干渉は絶対位相でなく、モード間の相対位相に依存する。

\hypertarget{ux5c40ux6240ux56deux8ee2ux5206ux5c90ux518dux7d50ux5408}{%
\section{局所回転、分岐、再結合}\label{ux5c40ux6240ux56deux8ee2ux5206ux5c90ux518dux7d50ux5408}}

実数 \(\epsilon_i(t)\) による局所位相蓄積は

\[
H_{\rm phase}
=
\sum_i
\epsilon_i(t)
\left|b_i\right|^2
\]

で実装できる。\(i\) と \(j\) の2モード混合は

\[
H_{ij}
=
g_{ij}(t)
\left(
b_i^*b_j+b_j^*b_i
\right)
\]

または位相をずらした Hermitian 結合で実装できる。有限時間の流れはユニタリ行列 \(U\) として

\[
b_{\rm out}
=
Ub_{\rm in},
\qquad
U^\dagger U=I
\]

と書ける。ここでユニタリ性は量子仮説でなく、全位相作用を保存する線形正準写像の複素表示である。

経路モード \(r\) が局所項 \(\epsilon_r(t)|b_r|^2\) に従うなら、相対位相は

\[
\theta_r(t)-\theta_s(t)
=
\theta_r(0)-\theta_s(0)
-
\frac1{\mathcal J_0}
\int_0^t
\left[
\epsilon_r(u)-\epsilon_s(u)
\right]
\,du
\]

となる。この式は位相担体が経路差の履歴を保持できることを示す。ただし、積分が粒子の古典作用差と一致することは別の導出を要する。

\hypertarget{ux5ba2ux89b3ux7684ux76f8ux95a2ux884cux5217}{%
\section{客観的相関行列}\label{ux5ba2ux89b3ux7684ux76f8ux95a2ux884cux5217}}

第1章の母測度 \(\mu_*\) を、共通調製条件 \(\mathcal P\) と測定プログラム \(M\) で条件付け、

\[
C_M(t)
=
\mathbb E_{\mu_*}
\left[
b_t b_t^\dagger
\mid
\mathcal P,M
\right]
\]

と定める。以下、条件を固定して混同がない場合は添字 \(M\) と時刻 \(t\) を省く。任意の \(v\in\mathbb C^L\) に対して

\[
v^\dagger Cv
=
\mathbb E_{\mu_*}
\left[
\left|v^\dagger b\right|^2
\mid
\mathcal P,M
\right]
\geq0
\]

なので \(C\) は正半定値 Hermitian 行列である。対角成分 \(C_{ii}\) は局所作用の集団平均、非対角成分 \(C_{ij}\) は相対位相相関を保持する。

\(C\) は調製条件とプログラムに依存するが、観測者の主観に依存するとは限らない。同じ源、Hamiltonian、母測度の条件付き集団を再現すれば同じ \(C\) が得られるという意味で、温度や流体密度と同様の客観的な集団状態量として扱う。\(C\) は独立の統計入力ではなく、同じ \(\mu_*\) の2次モーメントである。

\hypertarget{ux30bbux30ebux4f53ux7a4dux3092ux542bux3080ux898fux683cux5316}{%
\section{セル体積を含む規格化}\label{ux30bbux30ebux4f53ux7a4dux3092ux542bux3080ux898fux683cux5316}}

空間セルの体積を \(\Delta V\) とする。連続密度に対応する振幅 \(a_i\) を使う場合は、セル体積を吸収した正準振幅

\[
b_i
=
\sqrt{\Delta V}\,a_i
\]

を相関行列の定義に用いる。従って

\[
C_{ij}
=
\Delta V
\mathbb E_{\mu_*}
\left[
a_i a_j^*
\mid
\mathcal P,M
\right]
\]

であり、正規化対角成分

\[
p_i
=
\frac{C_{ii}}{\operatorname{tr}C}
\]

は既に規格化セル重みである。\(p_i\) へさらに \(\Delta V\) を掛けない。

階数1因子を \(\chi\) とし、\(\chi^\dagger\chi=1\) とする。連続密度表示を

\[
\psi_i
=
\frac{\chi_i}{\sqrt{\Delta V}}
\]

と定めれば、

\[
p_i
=
\left|\chi_i\right|^2
=
\left|\psi_i\right|^2
\Delta V,
\qquad
\sum_i p_i=1
\]

となる。

\hypertarget{ux968eux65701ux6761ux4ef6ux306eux610fux5473}{%
\section{階数1条件の意味}\label{ux968eux65701ux6761ux4ef6ux306eux610fux5473}}

\begin{theorem}[階数1相関と試行振幅の同値条件]
$C\neq0$ とする。次の2条件は同値である。

1. $C=\Lambda\chi\chi^\dagger$、$\Lambda>0$、$\chi^\dagger\chi=1$。
2. ある複素確率変数 $c^\omega$ が存在し、$b^\omega=c^\omega\chi$ がほとんど確実に成立する。

このとき $\Lambda=\mathbb E|c^\omega|^2$ である。
\end{theorem}

証明は付録Aに置く。この定理により、階数1は「平均すると1方向だけが残る」という弱い条件ではない。各試行の振幅ベクトルが共通射影方向 \(\chi\) にあることを要求する。絶対位相と全振幅は試行ごとに異なってよいが、相対振幅と相対位相は共通でなければならない。

従って、旧2成分場模型の標本間コヒーレント集中を単に弱い条件へ置き換えたとは言えない。準備問題の位置を、実在連続場の集中から有限共通源の出力方向へ移したのである。

\hypertarget{ux5171ux901aux6e90ux306bux3088ux308bux968eux65701ux6e96ux5099}{%
\section{共通源による階数1準備}\label{ux5171ux901aux6e90ux306bux3088ux308bux968eux65701ux6e96ux5099}}

入力モードを \(e_0\) とし、各試行で

\[
b_{\rm in}^\omega
=
e^{i\beta^\omega}e_0
\]

を準備する。\(\beta^\omega\) は試行ごとに任意でよい。有限正準準備回路 \(U_{\rm prep}\) が

\[
U_{\rm prep}e_0
=
\chi_0,
\qquad
\chi_0^\dagger\chi_0=1
\]

を満たすなら、

\[
b_{\rm out}^\omega
=
e^{i\beta^\omega}\chi_0
\]

となる。

\begin{proposition}[共通源による階数1相関]
上の準備では、$\mathbb E[b_{\rm out}]=0$ であっても

$$
C_{\rm out}
=
\chi_0\chi_0^\dagger
$$

が成立する。
\end{proposition}

共通絶対位相を固定する必要はない。必要なのは、1つの源から全モードの相対振幅と相対位相を同じ正準回路で作ることである。

\hypertarget{ux6e96ux5099ux56deux8defux306eux7269ux7406ux7684ux7bc4ux56f2}{%
\section{準備回路の物理的範囲}\label{ux6e96ux5099ux56deux8defux306eux7269ux7406ux7684ux7bc4ux56f2}}

任意の有限次元ユニタリ行列は、2モード混合と局所位相回転の有限列へ分解できる。従って、与えられた \(\chi_0\) を準備する有限2次 Hamiltonian 回路は構成できる。

ただし、この事実は任意の量子状態が自然に生成されることを意味しない。\(\chi_0\) を装置へ設計値として入れれば、その相対振幅を古典回路へ符号化しただけである。物理源、外部ポテンシャル、境界条件から特定の \(\chi_0\) が選ばれる機構は別に示す必要がある。

また、反復運転には次が必要である。

\begin{enumerate}
\def\labelenumi{\arabic{enumi}.}
\tightlist
\item
  源モードの作用を毎試行同じ範囲へ戻す。
\item
  前試行の相対位相情報を不要自由度へ移す。
\item
  外部記録を保ったまま有限装置を再初期化する。
\item
  源の失敗率を結果依存に除外しない。
\end{enumerate}

本稿は理想準備写像を明示するが、この全再初期化周期を完成していない。

\hypertarget{ux8fd1ux4f3cux968eux65701}{%
\section{近似階数1}\label{ux8fd1ux4f3cux968eux65701}}

\(C\) の固有値を \(\lambda_1\geq\lambda_2\geq\cdots\geq0\) とし、主固有ベクトルを \(\chi\) とする。階数欠陥を

\[
\varepsilon_{\rm rank}
=
1
-
\frac{\lambda_1}{\operatorname{tr}C}
\]

と定める。すると

\[
C
=
\lambda_1\chi\chi^\dagger
+
E,
\qquad
E\geq0,
\qquad
\operatorname{tr}E
=
\varepsilon_{\rm rank}
\operatorname{tr}C
\]

である。

純度欠陥

\[
\varepsilon_{\rm pur}
=
1
-
\frac{\operatorname{tr}C^2}{\left(\operatorname{tr}C\right)^2}
\]

も使えるが、節の残留強度には \(\varepsilon_{\rm rank}\) の方が直接的である。両者を同じ誤差として扱わない。

\hypertarget{ux9589ux9396ux7ddaux5f62ux767aux5c55ux3067ux306fux7d14ux5316ux3057ux306aux3044}{%
\section{閉鎖線形発展では純化しない}\label{ux9589ux9396ux7ddaux5f62ux767aux5c55ux3067ux306fux7d14ux5316ux3057ux306aux3044}}

第3章で示すように、閉鎖線形発展では

\[
C(t)
=
U(t)C(0)U(t)^\dagger
\]

となる。従って固有値と \(\varepsilon_{\rm rank}\) は保存される。

\begin{corollary}[閉鎖線形発展による純化の不可能性]
$\operatorname{rank}C(0)>1$ なら、有限時間の閉鎖線形 Hamiltonian 発展だけで $C(t)$ を階数1にできない。
\end{corollary}

これは現行模型の否定的結果である。階数1準備には共通源による初期化、不要成分の補助系への交換、条件付け、弱い外部交換の少なくとも1つが必要になる。後段の厳密な交換子発展をもって、準備まで導出したとは書かない。

\hypertarget{ux975eux7ddaux5f62ux9805ux3068ux9589ux5305ux6b8bux5dee}{%
\section{非線形項と閉包残差}\label{ux975eux7ddaux5f62ux9805ux3068ux9589ux5305ux6b8bux5dee}}

各試行が

\[
i\mathcal J_0\dot b^\omega
=
h(t)b^\omega
+
r^\omega
\]

に従う場合、相関行列は

\[
i\mathcal J_0\dot C
=
\left[h,C\right]
+
D_C,
\]

\[
D_C
=
\mathbb E_{\mu_*}
\left[
r b^\dagger
-
b r^\dagger
\mid
\mathcal P,M
\right]
\]

を満たす。Cauchy--Schwarz 不等式により

\[
\left\|D_C\right\|_{\rm op}
\leq
2
\left(
\mathbb E\left\|r\right\|^2
\right)^{1/2}
\left(
\mathbb E\left\|b\right\|^2
\right)^{1/2}
\]

である。

4次以上の Hamiltonian では \(r\) が \(b\) に非線形に依存し、\(D_C\) は一般に4次以上のモーメントを含む。\(C\) だけの閉包は自動的に成立しない。本稿は2次模型を厳密な基準とし、準2次模型では有限観測時間内の \(D_C\) を主要誤差として追跡する。

\hypertarget{ux968eux65701ux6761ux4ef6ux3068ux6a19ux672cux5316ux306eux5f79ux5272ux5206ux62c5}{%
\section{階数1条件と標本化の役割分担}\label{ux968eux65701ux6761ux4ef6ux3068ux6a19ux672cux5316ux306eux5f79ux5272ux5206ux62c5}}

階数1条件は、1つの統計振幅 \(\chi\) とその Schrödinger 型発展を取り出すために必要である。一方、第4章の正準作用選択器は各試行の実際の作用分配を読むため、有限基底 Born 型標本化それ自体には階数1を要求しない。

従って、次の2つを分ける。

\begin{enumerate}
\def\labelenumi{\arabic{enumi}.}
\tightlist
\item
  階数1の \(C\) から単一因子 \(\chi\) を得て、干渉振幅として伝播させること。
\item
  一般の正半定値 \(C\) に対し、固定全作用と条件付き一様角の下で有限基底結果を標本化すること。
\end{enumerate}

高階数相関行列は第1の目的には単一の \(\chi\) を与えないが、第2の目的では正規化行列 \(C/\operatorname{tr}C\) の対角要素を結果頻度として与え得る。

\hypertarget{ux76f8ux95a2ux884cux5217ux529bux5b66ux7a7aux9593ux30b0ux30e9ux30d5ux5e72ux6e09ux3068ux7bc0}{%
\chapter{相関行列力学、空間グラフ、干渉と節}\label{ux76f8ux95a2ux884cux5217ux529bux5b66ux7a7aux9593ux30b0ux30e9ux30d5ux5e72ux6e09ux3068ux7bc0}}

\begin{statusbox}
位置づけ：共通2次 Hamiltonian 下の相関行列発展、階数1因子の発展、有限グラフ上の干渉、階数欠陥による節の残留強度上界は厳密である。局所結合係数の必然的導出、連続極限、位相量子化、粒子の連続運動は未完成である。
\end{statusbox}

\hypertarget{ux76f8ux95a2ux884cux5217ux306eux9589ux3058ux305fux767aux5c55}{%
\section{相関行列の閉じた発展}\label{ux76f8ux95a2ux884cux5217ux306eux9589ux3058ux305fux767aux5c55}}

母測度を同じ調製条件 \(\mathcal P\) とプログラム \(M\) で条件付けた集団の各試行が、観測窓で共通の有限2次 Hamiltonian に従い、

\[
i\mathcal J_0\dot b^\omega
=
h(t)b^\omega
\]

を満たすとする。

\begin{theorem}[相関行列の交換子発展]
相関行列

$$
C_M(t)
=
\mathbb E_{\mu_*}
\left[
b^\omega(t)
\left(b^\omega(t)\right)^\dagger
\mid
\mathcal P,M
\right]
$$

は厳密に

$$
i\mathcal J_0\dot C_M
=
\left[h,C_M\right]
$$

を満たす。
\end{theorem}

\begin{proof}
$bb^\dagger$ を微分し、$i\mathcal J_0\dot b=hb$ とその共役転置を代入すると

$$
i\mathcal J_0
\frac{d}{dt}
\left(bb^\dagger\right)
=
hbb^\dagger
-
bb^\dagger h
$$

となる。集団平均を取る。
\end{proof}

この方程式は、単一の統計振幅 \(\chi\) を先に仮定しない。基礎的な集団状態は \(C_M\) である。同じ条件付き集団の中で \(h\) が試行、設定、ミクロ状態に依存する場合は平均と交換子を分離できず、単純な閉包は成立しない。準備、測定、再初期化を含む全周期にこの式を延長することも主張しない。

\hypertarget{ux4fddux5b58ux91cf}{%
\section{保存量}\label{ux4fddux5b58ux91cf}}

時間発展作用素を

\[
i\mathcal J_0\dot U(t,t_0)
=
h(t)U(t,t_0),
\qquad
U(t_0,t_0)=I
\]

とする。\(h=h^\dagger\) なら \(U\) はユニタリであり、

\[
C(t)
=
U(t,t_0)
C(t_0)
U(t,t_0)^\dagger
\]

となる。

\begin{corollary}[相関スペクトルの保存]
閉鎖線形発展では、$C$ の跡、全固有値、階数、および

$$
\mathcal P_C
=
\frac{\operatorname{tr}C^2}{\left(\operatorname{tr}C\right)^2}
$$

で定める純度が保存される。
\end{corollary}

従って、観測窓の理想2次発展は既に準備された相関を回転させるだけである。相関の純化または熱化を説明しない。

\hypertarget{ux968eux65701ux56e0ux5b50ux306e-schruxf6dinger-ux578bux767aux5c55}{%
\section{階数1因子の Schrödinger 型発展}\label{ux968eux65701ux56e0ux5b50ux306e-schruxf6dinger-ux578bux767aux5c55}}

\(C=\Lambda\chi\chi^\dagger\)、\(\chi^\dagger\chi=1\) とする。\(\Lambda=\operatorname{tr}C\) は保存される。

\begin{theorem}[階数1統計振幅の発展]
相関行列が階数1で交換子方程式に従うなら、ある実関数 $\lambda(t)$ が存在して

$$
i\mathcal J_0\dot\chi
=
h\chi
+
\lambda(t)\chi
$$

となる。時間依存共通位相を選べば、

$$
i\mathcal J_0\dot\chi
=
h\chi
$$

とできる。
\end{theorem}

証明は付録Aに置く。\(\lambda(t)\) は射影 \(\chi\chi^\dagger\) に影響しない共通位相だけを表す。

この定理は Schrödinger 型の有限次元形式を与えるが、\(h\) の物理的内容を自動的に決めない。任意の Hermitian 行列を選べること自体は、量子力学的構造の創発ではない。

\hypertarget{ux7a7aux9593ux30b0ux30e9ux30d5ux4e0aux306eux5c40ux6240-hamiltonian}{%
\section{空間グラフ上の局所 Hamiltonian}\label{ux7a7aux9593ux30b0ux30e9ux30d5ux4e0aux306eux5c40ux6240-hamiltonian}}

有限空間グラフ \(G=(V,E)\) を取り、各頂点に1つの位相担体を置く。辺係数 \(g_{ij}=g_{ji}\geq0\) は長さの逆2乗の次元を持つとする。グラフ Laplacian を

\[
\left(L_G\chi\right)_i
=
\sum_{j:\{i,j\}\in E}
g_{ij}
\left(
\chi_i-\chi_j
\right)
\]

と定める。局所外部ポテンシャルを \(V_i(t)\) とし、

\[
h_L(t)
=
\frac{\mathcal J_0^2}{2m}
L_G
+
\operatorname{diag}
\left(
V_1(t),\ldots,V_L(t)
\right)
\]

と置く。

対応する古典 Hamiltonian は

\[
H_{\rm ph}
=
\frac{\mathcal J_0^2}{2m}
\sum_{\{i,j\}\in E}
g_{ij}
\left|b_i-b_j\right|^2
+
\sum_i
V_i(t)
\left|b_i\right|^2
\]

である。各辺は隣接位相担体だけを結び、各ポテンシャル項は局所モードだけに作用する。\(Q\) と \(P\) の両直交成分へ同じ結合を置くため、共通内部回転が保たれる。

階数1因子は

\[
i\mathcal J_0\dot\chi
=
\left[
\frac{\mathcal J_0^2}{2m}L_G
+
V_L(t)
\right]
\chi
\]

に従う。規則格子で \(L_G\) が \(-\Delta\) の整合離散化なら、これは離散 Schrödinger 型方程式である。

\hypertarget{ux65e2ux77e5ux306eux53e4ux5178ux632fux52d5ux5b50ux8868ux793aux3068ux306eux533aux5225}{%
\section{既知の古典振動子表示との区別}\label{ux65e2ux77e5ux306eux53e4ux5178ux632fux52d5ux5b50ux8868ux793aux3068ux306eux533aux5225}}

量子状態を古典正準座標または結合振動子へ写し、有限次元 Schrödinger 方程式と同型の運動を得ること自体は既知である \cite{heslot1985,briggs_eisfeld2012,briggs_eisfeld2013,skinner2013}。従って、本稿は次を新規な導出として主張しない。

\begin{enumerate}
\def\labelenumi{\arabic{enumi}.}
\tightlist
\item
  複素ベクトルを2倍次元の実ベクトルで表すこと。
\item
  Hermitian 行列を2次古典 Hamiltonian として表すこと。
\item
  古典結合振動子がユニタリ行列と同型の線形発展を持つこと。
\end{enumerate}

本稿が検討する追加構造は、\(C\) を統一母測度の客観的集団相関として定義し、共通源による階数1準備、正準作用選択器、Bell設定条件付き履歴測度へ接続することである。

また、係数 \(\mathcal J_0^2/(2m)\) と局所ポテンシャル \(V_i\) は、本節では古典結合の設計値である。より原始的な粒子・媒質 Hamiltonian からこれらの値が必然的に現れることは未導出であり、Q1の未完成部分に残る。

\hypertarget{ux7d4cux8defux5e72ux6e09}{%
\section{2経路干渉}\label{ux7d4cux8defux5e72ux6e09}}

入力1モードを等分岐すると

\[
\chi_{\rm arm}
=
\frac1{\sqrt2}
\begin{pmatrix}
1\\
1
\end{pmatrix}
\]

となる。経路位相 \(\phi_1,\phi_2\) を蓄積し、同じ50対50混合で再結合すると

\[
\chi_+
=
\frac{
e^{i\phi_1}+e^{i\phi_2}
}{2},
\qquad
\chi_-
=
\frac{
e^{i\phi_1}-e^{i\phi_2}
}{2}
\]

となる。従って

\[
p_+
=
\cos^2
\left(
\frac{\phi_1-\phi_2}{2}
\right),
\qquad
p_-
=
\sin^2
\left(
\frac{\phi_1-\phi_2}{2}
\right)
\]

である。\(\phi_1-\phi_2=\pi\) では \(p_+=0\) となり、出力統計に完全な節が生じる。

この干渉は正の対角成分の混合だけでは生じない。再結合前の非対角相関

\[
C_{12}
=
\Lambda
\chi_1\chi_2^*
\]

が相対位相を保持することが必要である。

\hypertarget{ux968eux6570ux6b20ux9665ux3068ux7bc0ux306eux6df1ux3055}{%
\section{階数欠陥と節の深さ}\label{ux968eux6570ux6b20ux9665ux3068ux7bc0ux306eux6df1ux3055}}

\(C=\lambda_1\chi\chi^\dagger+E\)、\(E\geq0\) とし、ユニタリ再結合 \(U\) の出力 \(k\) が理想因子に対して節を持つ、すなわち \((U\chi)_k=0\) とする。

\begin{proposition}[階数欠陥による節の残留強度上界]
正規化出力強度

$$
p_k
=
\frac{
\left(UCU^\dagger\right)_{kk}
}{
\operatorname{tr}C
}
$$

は

$$
p_k
\leq
\varepsilon_{\rm rank}
$$

を満たす。
\end{proposition}

\begin{proof}
節条件により主成分の寄与は零であり、

$$
p_k
=
\frac{
\left(UEU^\dagger\right)_{kk}
}{
\operatorname{tr}C
}
\leq
\frac{\left\|E\right\|_{\rm op}}{\operatorname{tr}C}
\leq
\frac{\operatorname{tr}E}{\operatorname{tr}C}
=
\varepsilon_{\rm rank}
$$

である。
\end{proof}

従って、節の残留強度は階数1からのずれを直接検査する量になる。

\hypertarget{ux4f4dux76f8ux96d1ux97f3ux3068ux53efux8996ux5ea6}{%
\section{位相雑音と可視度}\label{ux4f4dux76f8ux96d1ux97f3ux3068ux53efux8996ux5ea6}}

相対位相ゆらぎを \(\delta\phi\) とする。平均干渉項は

\[
\mathbb E
\left[
e^{i\delta\phi}
\right]
\]

だけ減衰する。零平均 Gauss ゆらぎで分散が \(\sigma_\phi^2\) なら、可視度は

\[
V_{\rm int}
=
e^{-\sigma_\phi^2/2}
\]

である。小さい位相雑音では \(1-V_{\rm int}=O(\sigma_\phi^2)\) となる。階数欠陥と位相雑音は異なる機構なので、検算では別に測る。

\hypertarget{ux9023ux7d9aux88dcux9593ux3068ux96e2ux6563ux5316ux8aa4ux5dee}{%
\section{連続補間と離散化誤差}\label{ux9023ux7d9aux88dcux9593ux3068ux96e2ux6563ux5316ux8aa4ux5dee}}

規則格子のセル中心 \(x_i\) に \(\psi_i=\chi_i/\sqrt{\Delta V}\) を置き、補間作用素を \(\mathcal I_L\) とする。滑らかな試験関数に対する整合性を

\[
\left\|
L_G\mathcal R_L f
+
\mathcal R_L\Delta f
\right\|
\leq
c_L\ell^p
\left\|f\right\|_{H^{p+2}}
\]

と書く。\(\mathcal R_L\) はセル標本化、\(\ell\) は格子幅である。

有限時間 \(0\leq t\leq T\) で安定性があれば、連続 Schrödinger 型方程式との差は概ね \(O(T\ell^p)\) で蓄積する。境界、非一様格子、時間依存ポテンシャルでは別の安定性定数が必要である。本稿は有限グラフの厳密結果を中心とし、一般連続極限の一様定理を主張しない。

\hypertarget{ux56faux6709ux30d9ux30afux30c8ux30ebux306eux6761ux4ef6ux4ed8ux304dux5305ux542b}{%
\section{固有ベクトルの条件付き包含}\label{ux56faux6709ux30d9ux30afux30c8ux30ebux306eux6761ux4ef6ux4ed8ux304dux5305ux542b}}

時間非依存 \(h_L\) の固有ベクトル \(\varphi_n\) が

\[
h_L\varphi_n
=
E_n\varphi_n
\]

を満たすなら、

\[
\chi_n(t)
=
e^{-iE_nt/\mathcal J_0}
\varphi_n
\]

は厳密解である。従って、離散ポテンシャル模型は節を持つ固有ベクトルを解として含む。

これは次を意味しない。

\begin{enumerate}
\def\labelenumi{\arabic{enumi}.}
\tightlist
\item
  \(h_L\) のスペクトルが実在粒子のエネルギー測定値になること。
\item
  任意初期集団が特定の固有ベクトルへ吸引されること。
\item
  基底状態または励起状態が弱開放系で選択されること。
\item
  井戸型または調和型ポテンシャルの連続極限が定量的に再現されたこと。
\end{enumerate}

\hypertarget{ux7c92ux5b50ux306eux9023ux7d9aux904bux52d5ux3092ux5c0eux304bux306aux3044}{%
\section{粒子の連続運動を導かない}\label{ux7c92ux5b50ux306eux9023ux7d9aux904bux52d5ux3092ux5c0eux304bux306aux3044}}

連続密度表示から形式的な統計流

\[
j_\psi
=
\frac{\mathcal J_0}{m}
\operatorname{Im}
\left(
\psi^*\nabla\psi
\right)
\]

を作れる。しかし \(C\) と \(\psi\) は集団量であり、単一試行の粒子 Hamiltonian に \(j_\psi/|\psi|^2\) を直接入れると、求める集団量をミクロ法則へ戻す循環になる。

必要なのは、単一試行の変数だけからなる

\[
\dot X^\omega
=
F_X
\left(
X^\omega,
P_X^\omega,
b^\omega,
Y^\omega
\right)
\]

を先に構成し、その集団流束が \(j_\psi\) へ縮約されることを示すことである。本稿はこの流束定理を持たない。

従って、主張は再結合後の統計強度と第4章の有限基底位相担体標本化までに限定する。2重スリットの粒子検出、途中の粒子軌道、節を避ける運動、全時刻での等変性は未解決である。

\hypertarget{ux4f4dux76f8ux91cfux5b50ux5316ux306eux4f4dux7f6eux3065ux3051}{%
\section{位相量子化の位置づけ}\label{ux4f4dux76f8ux91cfux5b50ux5316ux306eux4f4dux7f6eux3065ux3051}}

旧2成分誘導場では、非零の単価な連続場に対する閉路巻数から条件付き循環量子化を記述できた。現行模型は有限グラフ上の位相担体を基礎とし、一般の閉路で位相差が定義できても、連続空間の単価性、節の生成・消滅、閉路変形に対する巻数保存をまだ統合していない。

従って、Wallstrom 問題は未解決へ戻す。将来の課題は、非零位相担体が置かれた空間グラフ、辺位相差、閉路巻数、連続補間、節を通る位相すべりを同じ有限 Hamiltonian で記述することである。

\hypertarget{ux6b63ux6e96ux4f5cux7528ux9078ux629eux5668ux3068ux6709ux9650ux57faux5e95ux6a19ux672cux5316}{%
\chapter{正準作用選択器と有限基底標本化}\label{ux6b63ux6e96ux4f5cux7528ux9078ux629eux5668ux3068ux6709ux9650ux57faux5e95ux6a19ux672cux5316}}

\begin{statusbox}
位置づけ：条件付き一様な正準角から、各試行で作用比に一致する排他的な1結果を得る式は厳密である。固定全作用と基底に依存しない調製相関の下で、一般の正半定値相関行列に対する有限基底 Born 型則を得る。滑らかな比較器、粒子検出、完全周期は未完成である。
\end{statusbox}

\hypertarget{ux6e2cux5b9aux57faux5e95ux3068ux51faux529bux4f5cux7528}{%
\section{測定基底と出力作用}\label{ux6e2cux5b9aux57faux5e95ux3068ux51faux529bux4f5cux7528}}

\(L\) 個の位相担体モードを \(b\in\mathbb C^L\) とする。測定する有限直交基底をユニタリ行列 \(U\) で表し、測定前の正準モード混合を

\[
b'
=
Ub
\]

とする。複素ユニタリ変換は実正準変数 \((Q,P)\) 上の直交正準変換に対応する。出力モード \(k\) の作用と全位相作用を

\[
I_k
=
\mathcal J_0
\left|b'_k\right|^2,
\qquad
I_{\rm ph}
=
\sum_{k=1}^L I_k
=
\mathcal J_0b^\dagger b
\]

とする。\(U\) は全作用を保存する。以下、\(I_{\rm ph}>0\) がほとんど確実に成立する条件付き集団を扱う。

\(I_k\) は各試行に実在する正準作用である。装置は集団量 \(C\) を単一試行の Hamiltonian へ代入せず、各試行の \(I_k\) だけを読み出す。

\hypertarget{ux56faux5b9aux4f5cux7528ux306eux6b63ux6e96ux9078ux629eux5668}{%
\section{固定作用の正準選択器}\label{ux56faux5b9aux4f5cux7528ux306eux6b63ux6e96ux9078ux629eux5668}}

補助振動子の作用角変数を

\[
\left(
J_{\rm sel},
\vartheta_{\rm sel}
\right),
\qquad
J_{\rm sel}=J_*>0
\]

とする。選択器角に必要なのは周辺一様性ではなく、位相担体、測定基底、調製条件の下での条件付き一様性である。具体的に、

\[
\mu_*
\left(
d\vartheta_{\rm sel}
\mid
b,U,\mathcal P,M=\mathsf{basis}
\right)
=
\frac{d\vartheta_{\rm sel}}{2\pi}
\]

を仮定する。\(\vartheta_{\rm sel}\) の周辺分布だけが一様でも、作用分配と相関していれば以下の区間則は一般に成立しない。

選択変数を

\[
u
=
\frac{\vartheta_{\rm sel}}{2\pi}
I_{\rm ph},
\qquad
0\leq u<I_{\rm ph}
\]

とする。角度の切断点 \(0=2\pi\) は零測度境界であり、有限幅実装では比較誤差として扱う。

\hypertarget{ux7d2fux7a4dux4f5cux7528ux533aux9593ux3068ux552fux4e00ux7d50ux679c}{%
\section{累積作用区間と唯一結果}\label{ux7d2fux7a4dux4f5cux7528ux533aux9593ux3068ux552fux4e00ux7d50ux679c}}

累積作用を

\[
S_0=0,
\qquad
S_k
=
\sum_{j=1}^k I_j
\]

とする。結果 \(k\) の事象を

\[
E_k^U
=
\left\{
S_{k-1}\leq u<S_k
\right\}
\]

と定める。\(S_L=I_{\rm ph}\) なので、境界を除いて事象 \(E_k^U\) は互いに素であり、全標本を覆う。従って、各試行はちょうど1つの結果を持つ。

この結果は粒子が入口 \(k\) を最初に通過した事象ではない。結果レジスターが位相担体の出力モード \(k\) を記録した事象である。旧位置作用殻の共通入口流束、全入口の排他的作用転送、2モード殻容量は本定理の仮定から外れる。

\hypertarget{ux4e3bux5b9aux74062ux306eux4e00ux822cux5f62}{%
\section{主定理2の一般形}\label{ux4e3bux5b9aux74062ux306eux4e00ux822cux5f62}}

\begin{theorem}[正準作用選択器による有限基底標本化]
次を仮定する。

1. $I_{\rm ph}>0$ がほとんど確実に成立する。
2. 測定前の正準混合 $U$ が全作用を保存する。
3. 選択器角が $(b,U,\mathcal P)$ の下で条件付き一様である。
4. 累積作用区間が排他的に比較され、各試行で1つの結果だけが記録される。
5. 比較失敗、記録失敗、再帰失敗を結果に応じて捨てない。

このとき各試行の条件付き結果確率は

$$
P(k\mid b,U)
=
\frac{I_k}{I_{\rm ph}}
=
\frac{
\left|\left(Ub\right)_k\right|^2
}{
b^\dagger b
}
$$

であり、母測度で平均した頻度は

$$
P(k\mid U,\mathcal P)
=
\mathbb E_{\mu_*}
\left[
\frac{I_k}{I_{\rm ph}}
\middle|
U,\mathcal P,M=\mathsf{basis}
\right]
$$

である。
\end{theorem}

\begin{proof}
$(b,U)$ を固定すると、$u$ は区間 $[0,I_{\rm ph})$ 上で一様である。$E_k^U$ の区間長は $S_k-S_{k-1}=I_k$ なので、条件付き確率は $I_k/I_{\rm ph}$ である。条件付き期待値を取ると第2式を得る。
\end{proof}

周期写像が \(\mu_*\) に関して不変かつエルゴード的であるという第1章の仮定により、同じ装置を繰り返した結果 \(k\) の長期相対頻度は、この母測度確率へ収束する。

\hypertarget{ux56faux5b9aux5168ux4f5cux7528ux3068-born-ux578bux5247}{%
\section{固定全作用と Born 型則}\label{ux56faux5b9aux5168ux4f5cux7528ux3068-born-ux578bux5247}}

測定基底 \(U\) の下の相関行列を

\[
C_U
=
\mathbb E_{\mu_*}
\left[
bb^\dagger
\mid
U,\mathcal P,M=\mathsf{basis}
\right]
\]

とする。全位相作用が条件付き集団で固定され、

\[
I_{\rm ph}=I_0>0
\]

なら、

\[
\operatorname{tr}C_U
=
\frac{I_0}{\mathcal J_0}
\]

である。

\begin{corollary}[固定作用下の有限基底 Born 型則]
主定理2の条件に加えて $I_{\rm ph}=I_0$ を仮定すると、

$$
P(k\mid U,\mathcal P)
=
\frac{
\left(
UC_UU^\dagger
\right)_{kk}
}{
\operatorname{tr}C_U
}
$$

が成立する。
\end{corollary}

階数1は仮定していない。従って、純度が1未満の高階数相関行列も同じ選択器で扱える。

\hypertarget{ux5168ux4f5cux7528ux5909ux52d5ux306eux5171ux5206ux6563ux88dcux6b63}{%
\section{全作用変動の共分散補正}\label{ux5168ux4f5cux7528ux5909ux52d5ux306eux5171ux5206ux6563ux88dcux6b63}}

一般には比の平均と平均の比は一致しない。正規化作用分配を

\[
r_k
=
\frac{I_k}{I_{\rm ph}}
\]

とすると、主定理2は \(P_k=\mathbb E[r_k]\) を与える。一方、

\[
\frac{
\mathbb E[I_k]
}{
\mathbb E[I_{\rm ph}]
}
=
\mathbb E[r_k]
+
\frac{
\operatorname{Cov}
\left(
I_{\rm ph},r_k
\right)
}{
\mathbb E[I_{\rm ph}]
}
\]

なので、

\[
P_k
-
\frac{
\left(
UC_UU^\dagger
\right)_{kk}
}{
\operatorname{tr}C_U
}
=
-
\frac{
\operatorname{Cov}
\left(
I_{\rm ph},r_k
\right)
}{
\mathbb E[I_{\rm ph}]
}
\]

である。従って、Born 型の平均比と一致する十分条件は次のいずれかである。

\begin{enumerate}
\def\labelenumi{\arabic{enumi}.}
\tightlist
\item
  \(I_{\rm ph}\) が試行ごとに固定される。
\item
  \(I_{\rm ph}\) と \(r_k\) が無相関である。
\end{enumerate}

全作用変動を許すとき、この共分散項を消えたものとして扱わない。

\hypertarget{ux4efbux610fux6709ux9650ux57faux5e95ux3068ux8abfux88fdux72ecux7acbux6027}{%
\section{任意有限基底と調製独立性}\label{ux4efbux610fux6709ux9650ux57faux5e95ux3068ux8abfux88fdux72ecux7acbux6027}}

同じ状態 \(C\) を異なる基底 \(U\) で測ったと主張するには、単粒子側の調製相関が基底選択に依存しないことが必要である。すなわち、

\[
C_U=C
\]

を要求する。固定作用下では

\[
P(k\mid U,\mathcal P)
=
\frac{
\left(
UCU^\dagger
\right)_{kk}
}{
\operatorname{tr}C
}
\]

となる。これは有限次元の任意の直交基底に対する Born 型則である。

Bell側では設定条件付き母測度を許すが、その測定設定依存性を単粒子側へ無条件に持ち込まない。\(C_U\neq C\) なら、上式は同じ状態の基底変更でなく、基底に応じて異なる条件付き集団を測っている。

\hypertarget{ux7a7aux9593ux30bbux30ebux57faux5e95}{%
\section{空間セル基底}\label{ux7a7aux9593ux30bbux30ebux57faux5e95}}

\(U=I\) とし、モード \(i\) が体積 \(\Delta V\) の空間セルに対応するとする。第2章の規格化により、固定作用下では

\[
P_i
=
\frac{C_{ii}}{\operatorname{tr}C}
\]

である。階数1なら

\[
C
=
\Lambda\chi\chi^\dagger,
\qquad
\psi_i
=
\frac{\chi_i}{\sqrt{\Delta V}}
\]

なので、

\[
P_i
=
\left|\chi_i\right|^2
=
\left|\psi_i\right|^2
\Delta V
\]

となる。

ただし、これは位相担体の空間セル基底を標本化する式である。粒子座標 \(X\) がセル \(i\) の局所検出器を作動させたこと、粒子と選択モードが同期すること、粒子が連続軌道を持つことは別の導出を要する。従って本稿では「粒子位置 Born 則」ではなく「空間セル基底の Born 型標本化」と呼ぶ。

\hypertarget{ux53efux9006ux88c5ux7f6eux306eux5019ux88dc}{%
\section{可逆装置の候補}\label{ux53efux9006ux88c5ux7f6eux306eux5019ux88dc}}

理想装置は次の順序で働く。

\begin{enumerate}
\def\labelenumi{\arabic{enumi}.}
\tightlist
\item
  正準回路 \(U\) が測定基底へモードを混合する。
\item
  各 \(I_k\) を累積作用レジスターへ可逆に加える。
\item
  選択器角から \(u\) を計算する。
\item
  \(u-S_k\) の符号変化を滑らかな有限幅比較器で検出する。
\item
  最初の1結果だけをラッチへ記録する。
\item
  結果を空の外部記録へ可逆にコピーする。
\item
  比較、累積加算、基底混合を逆計算する。
\item
  消せない結果情報と不要情報を外部自由度へ移す。
\end{enumerate}

正準加算と比較器の候補 Hamiltonian は付録Cに示す。直列走査の相互作用窓は \(O(L)\) 個であり、二分木比較では深さを \(O(\log L)\) にできるが、追加レジスターが必要である。本稿は構成要素を示すが、全てを1つの滑らかな自律 Hamiltonian と不変母測度へ統合していない。

\hypertarget{ux6761ux4ef6ux4ed8ux304dux4e00ux69d8ux6027ux3068ux6bd4ux8f03ux5e45ux306eux8aa4ux5dee}{%
\section{条件付き一様性と比較幅の誤差}\label{ux6761ux4ef6ux4ed8ux304dux4e00ux69d8ux6027ux3068ux6bd4ux8f03ux5e45ux306eux8aa4ux5dee}}

選択器角の条件付き分布と一様分布の全変動距離を

\[
\varepsilon_{\rm sel}
=
\mathbb E
\left[
d_{\rm TV}
\left(
\mu_*
\left(
d\vartheta_{\rm sel}\mid b,U,\mathcal P,M=\mathsf{basis}
\right),
\frac{d\vartheta_{\rm sel}}{2\pi}
\right)
\right]
\]

とする。このとき各結果確率の一様選択器からの偏差は \(\varepsilon_{\rm sel}\) 以下である。

比較境界の半幅を \(w>0\) とし、

\[
\varepsilon_{\rm cmp}(w)
=
\mu_*
\left(
\min_k
\left|u-S_k\right|
\leq w
\middle|
\mathcal P,M=\mathsf{basis}
\right)
\]

とする。境界近傍だけで理想結果と有限幅結果が異なる装置なら、比較器による全変動誤差は \(\varepsilon_{\rm cmp}(w)\) で抑えられる。理想境界が零測度で、条件付き密度が有界なら \(w\to0\) でこの誤差は零へ近づく。

\hypertarget{ux4e3bux5f35ux306eux7bc4ux56f2}{%
\section{主張の範囲}\label{ux4e3bux5f35ux306eux7bc4ux56f2}}

本章で得たものは次である。

\begin{enumerate}
\def\labelenumi{\arabic{enumi}.}
\tightlist
\item
  各試行の作用分配から1結果を排他的に選ぶ厳密な区間則。
\item
  高階数を含む一般の正半定値相関行列への拡張。
\item
  固定作用下の任意有限直交基底 Born 型則。
\item
  全作用変動に対する厳密な共分散補正。
\item
  条件付き一様性と有限比較幅の誤差指標。
\end{enumerate}

一方、次は示していない。

\begin{enumerate}
\def\labelenumi{\arabic{enumi}.}
\tightlist
\item
  選択器角の条件付き一様性を具体的な周期写像から導くこと。
\item
  累積、比較、ラッチ、記録、逆計算を統合した完成 Hamiltonian。
\item
  連続スペクトルと無限次元測定。
\item
  位相担体モードの選択と粒子の局所位置検出の同一性。
\item
  粒子の連続軌道。
\end{enumerate}

\hypertarget{bell-ux4f59ux5f26ux679dux4f5cux7528ux3068ux8a2dux5b9aux6761ux4ef6ux4ed8ux304dux5c65ux6b74ux6e2cux5ea6}{%
\chapter{Bell 余弦枝作用と設定条件付き履歴測度}\label{bell-ux4f59ux5f26ux679dux4f5cux7528ux3068ux8a2dux5b9aux6761ux4ef6ux4ed8ux304dux5c65ux6b74ux6e2cux5ea6}}

\begin{statusbox}
位置づけ：反対称左右交差相関と局所回転からの余弦枝作用は補助模型内部で厳密である。Bell共同確率は、同じ不変母測度の設定条件付き完結履歴として定義する。余弦頻度には二側履歴整合、条件付き一様な境界角、4結果で共通な基準セクター密度が必要であり、その物理的生成は未完成である。
\end{statusbox}

\hypertarget{ux56faux5b9aux3055ux308cux305fux5de6ux53f3ux4f4dux76f8ux62c5ux4f53ux30e2ux30fcux30c9}{%
\section{固定された左右位相担体モード}\label{ux56faux5b9aux3055ux308cux305fux5de6ux53f3ux4f4dux76f8ux62c5ux4f53ux30e2ux30fcux30c9}}

Bell 構成では、左右へ伝播する有限個の局在位相担体モードを基礎変数とする。2粒子配置空間上の場も、単粒子側の統計振幅 \(\chi\) も局所 Hamiltonian へ追加しない。装置の組立時に固定した左右モードを

\[
z^A_{\mu r},
\qquad
z^B_{\nu r},
\qquad
\mu,\nu\in\{+,-\},
\qquad
r\in\{1,2\}
\]

とする。\(\mu,\nu\) は局所分析器の2出力、\(r\) は2つの直交源チャンネルである。設定または結果ごとにモード基底を選び直さない。

左右モードの派生交差相関を

\[
\Xi_{\mu\nu}
=
\sum_{r=1}^2
\eta_r
z^A_{\mu r}
\left(
z^B_{\nu r}
\right)^*
\]

と定める。\(\Xi\) は独立した場でも新しい正準変数でもない。左右の基礎位相担体から計算される階数2以下の交差相関である。単粒子側の \(C\) は正半定値 Hermitian 行列であるのに対し、\(\Xi\) は一般に正半定値でも Hermitian でもない。両者を同じ相関行列として扱わない。

\hypertarget{ux53cdux5bfeux79f0ux6e90ux3068ux5c40ux6240ux5206ux6790ux5668}{%
\section{反対称源と局所分析器}\label{ux53cdux5bfeux79f0ux6e90ux3068ux5c40ux6240ux5206ux6790ux5668}}

理想源が準備する交差相関を

\[
\Xi_0
=
\sqrt{
\frac{\mathcal K}{2}
}
\begin{pmatrix}
0&1\\
-1&0
\end{pmatrix},
\qquad
\mathcal K>0
\]

とする。これは2つの階数1交差項へ分解できるため、2つの直交源チャンネルで構成できる。有限 Hamiltonian 源が任意初期状態から \(\Xi_0\) を反復準備することと、比較時刻まで相対位相を保つことは独立の準備問題である。

Bell設定を \(x,y\)、対応する局所モード角を \(\alpha_x,\beta_y\) とする。左右の局所分析器は各側の基礎位相担体だけを

\[
z^{A,(r)}
\longmapsto
R(\alpha_x)
z^{A,(r)},
\qquad
z^{B,(r)}
\longmapsto
R(\beta_y)
z^{B,(r)}
\]

と回転する。ここで

\[
R(\theta)
=
\begin{pmatrix}
\cos\theta&-\sin\theta\\
\sin\theta&\cos\theta
\end{pmatrix}
\]

である。交差相関表示は

\[
\Xi(x,y)
=
R(\alpha_x)
\Xi_0
R(\beta_y)^{\mathsf T}
\]

となるが、物理的に回すのは空間的に分離した左右の局所モードである。

\hypertarget{ux4f59ux5f26ux679dux4f5cux7528}{%
\section{余弦枝作用}\label{ux4f59ux5f26ux679dux4f5cux7528}}

結果 \(A,B\in\{+1,-1\}\) に対応する基底を \(e_A,e_B\) とし、

\[
\Xi_{AB}^{xy}
=
e_A^{\mathsf T}
\Xi(x,y)
e_B,
\qquad
K_{AB}^{xy}
=
\left|
\Xi_{AB}^{xy}
\right|^2
\]

と定める。

\begin{theorem}[反対称対モードの Bell 余弦枝作用]
反対称源と左右の実回転の下で、

$$
K_{AB}^{xy}
=
\frac{\mathcal K}{4}
\left[
1
-
AB\cos\Delta_{xy}
\right],
\qquad
\Delta_{xy}
=
2
\left(
\alpha_x-\beta_y
\right)
$$

が成立する。さらに、

$$
\sum_{A,B}
K_{AB}^{xy}
=
\mathcal K,
\qquad
\sum_B
K_{AB}^{xy}
=
\sum_A
K_{AB}^{xy}
=
\frac{\mathcal K}{2}
$$

である。
\end{theorem}

\begin{proof}
$\delta=\alpha_x-\beta_y$ とする。行列積から

$$
K_{++}^{xy}
=
K_{--}^{xy}
=
\frac{\mathcal K}{2}
\sin^2\delta,
\qquad
K_{+-}^{xy}
=
K_{-+}^{xy}
=
\frac{\mathcal K}{2}
\cos^2\delta
$$

を得る。$\cos2\delta=\cos\Delta_{xy}$ を用いる。
\end{proof}

この定理は設定依存の非負枝作用を与えるが、結果頻度をまだ与えない。

\hypertarget{ux540cux3058ux6bcdux6e2cux5ea6ux4e0aux306eux5b8cux5168ux5c65ux6b74ux4e8bux8c61}{%
\section{同じ母測度上の完全履歴事象}\label{ux540cux3058ux6bcdux6e2cux5ea6ux4e0aux306eux5b8cux5168ux5c65ux6b74ux4e8bux8c61}}

第1章の母測度 \(\mu_*\) は、設定生成器、局所装置、比較器、記録、外部自由度まで含む開始面 \(\Sigma_0\) 上に置かれる。Bellプログラムの完結履歴事象を

\[
F_{AB}^{xy}
=
\left\{
z\in\Sigma_0:
x(z)=x,
y(z)=y,
A(z)=A,
B(z)=B,
M(z)=\mathsf{Bell}
\right\}
\]

とする。Bell共同確率は新しいBell専用測度でなく、

\[
P(A,B\mid x,y)
=
\mu_*
\left(
F_{AB}^{xy}
\mid
x,y,M=\mathsf{Bell}
\right)
\]

である。第1章で仮定した不変性とエルゴード性により、設定対 \((x,y)\) が出現する部分列の長期相対頻度をこの条件付き確率と結び付ける。各設定対の母測度が正であることを仮定する。

\hypertarget{ux5c40ux6240ux8a18ux9332ux3068ux5171ux901aux672aux6765}{%
\section{局所記録と共通未来}\label{ux5c40ux6240ux8a18ux9332ux3068ux5171ux901aux672aux6765}}

本稿の局所装置は、既に形成された2値結果セクターを出力ポートと指針へ写す最小符号化器である。連続したミクロ状態から安定な唯一結果を形成する一般測定器は構成していない。

左右の局所結果は、戻り信号が共通未来へ到達する前に記録される。理想局所応答は

\[
A
=
A(x,\Lambda_A),
\qquad
B
=
B(y,\Lambda_B)
\]

と書ける。ここで \(\Lambda_A,\Lambda_B\) は各側へ到達するミクロ変数である。後の比較器は過去の記録を変更しない。

従って、共通未来で \(K_{AB}^{xy}\) を計算してから \((A,B)\) を前向きに選ぶ模型は採用しない。共通未来の役割は、既に局所記録された結果と未来境界変数が整合する完結履歴を定めることである。

\hypertarget{bell-ux5883ux754cux89d2ux3068ux6574ux5408ux4e8bux8c61}{%
\section{Bell 境界角と整合事象}\label{bell-ux5883ux754cux89d2ux3068ux6574ux5408ux4e8bux8c61}}

共通未来に固定作用の選択器振動子を置き、その角を \(\vartheta_{\rm B}\) とする。枝作用の総和は \(\mathcal K\) なので、

\[
u_{\rm B}
=
\frac{\vartheta_{\rm B}}{2\pi}
\mathcal K,
\qquad
0\leq u_{\rm B}<\mathcal K
\]

とする。4結果を固定した順に並べ、累積枝作用を

\[
T_0^{xy}=0,
\qquad
T_m^{xy}
=
\sum_{n=1}^m
K_n^{xy}
\]

とする。結果対 \((A,B)\) に対応する区間 \(\mathcal I_{AB}^{xy}\) の長さは

\[
\left|
\mathcal I_{AB}^{xy}
\right|
=
K_{AB}^{xy}
\]

である。

完結履歴の整合事象を

\[
G
=
\left\{
u_{\rm B}
\in
\mathcal I_{A(z)B(z)}^{x(z)y(z)}
\right\}
\]

とする。物理的母測度には

\[
\mu_*
\left(
G
\mid
M=\mathsf{Bell}
\right)
=
1
\]

を要求する。実験で \(G\) を満たさない周期を後から捨てるなら事後選別であり、本稿の模型ではない。必要なのは、不整合周期が物理的解空間で零測度になることである。

\hypertarget{ux57faux6e96ux30bbux30afux30bfux30fcux5bc6ux5ea6}{%
\section{基準セクター密度}\label{ux57faux6e96ux30bbux30afux30bfux30fcux5bc6ux5ea6}}

整合条件を課す前の候補完結履歴に対する Liouville 基準要素を \(\nu_*^{xy}\) とする。これは別の観測統計ではなく、\(\mu_*\) のBell条件付き部分を構成するための境界位相体積である。結果セクター \((A,B)\) で、選択器座標以外を積分した基準密度を \(q_{AB}^{xy}>0\) とする。

選択器角が候補履歴の残る変数に対して条件付き一様で、整合条件以外の因子を \(q_{AB}^{xy}\) へ含めると、整合履歴の未規格化質量は

\[
W_{AB}^{xy}
=
q_{AB}^{xy}
K_{AB}^{xy}
\]

となる。従って一般式は

\[
P(A,B\mid x,y,G)
=
\frac{
q_{AB}^{xy}
K_{AB}^{xy}
}{
\sum_{A',B'}
q_{A'B'}^{xy}
K_{A'B'}^{xy}
}
\]

である。

\(q_{AB}^{xy}\) は、局所結果形成後のセクター密度、時計流束、余面積 Jacobian、付随自由度の体積、解多重度、有限境界幅を含む。正準角を追加しただけでは \(q_{AB}^{xy}\) の共通性は従わない。これは旧共通境界測度の課題を別名で消したものではなく、その残存条件を1つの因子として明示した式である。

\hypertarget{ux4e3bux5b9aux74063}{%
\section{主定理3}\label{ux4e3bux5b9aux74063}}

\begin{theorem}[共通基準密度下の Bell 完全履歴統計]
次を仮定する。

1. 第5.3節の余弦枝作用が成立する。
2. 局所結果が共通未来より前に形成され、左右の局所応答が保たれる。
3. Bell境界角が候補履歴の残る変数の下で条件付き一様である。
4. 物理的母測度が整合事象 $G$ に支持を持つ。
5. 各設定対で $q_{AB}^{xy}=q^{xy}>0$ が4結果に共通である。
6. 結果依存の失敗周期を捨てない。

このとき、同じ母測度の設定条件付き共同確率は

$$
P(A,B\mid x,y)
=
\frac14
\left[
1
-
AB\cos\Delta_{xy}
\right]
$$

となる。
\end{theorem}

\begin{proof}
共通 $q^{xy}$ は規格化で相殺する。第5.3節の $\sum_{A,B}K_{AB}^{xy}=\mathcal K$ を用いれば、$P(A,B\mid x,y)=K_{AB}^{xy}/\mathcal K$ となる。
\end{proof}

定理は、整合支持条件と共通基準密度を置いた後の厳密結果である。具体的な周期写像がこの母測度を生成することは示していない。

\hypertarget{ux975eux4fe1ux53f7ux5468ux8fbaux3068-chsh-ux5024}{%
\section{非信号周辺と CHSH 値}\label{ux975eux4fe1ux53f7ux5468ux8fbaux3068-chsh-ux5024}}

\begin{corollary}[対称履歴測度下の非信号性]
主定理3の条件の下で、

$$
\sum_B
P(A,B\mid x,y)
=
\frac12,
\qquad
\sum_A
P(A,B\mid x,y)
=
\frac12
$$

である。
\end{corollary}

共同相関は

\[
E(x,y)
=
\sum_{A,B}
AB
P(A,B\mid x,y)
=
-\cos\Delta_{xy}
\]

となる。標準4設定では

\[
\left|
S_{\rm CHSH}
\right|
=
2\sqrt2
\]

を得る。これは平面内2出力の理想余弦則に対する値であり、一般 Tsirelson 原理の導出ではない。

\hypertarget{ux6e2cux5b9aux8a2dux5b9aux72ecux7acbux6027ux3068ux8a2dux5b9aux4ecbux5165}{%
\section{測定設定独立性と設定介入}\label{ux6e2cux5b9aux8a2dux5b9aux72ecux7acbux6027ux3068ux8a2dux5b9aux4ecbux5165}}

設定生成器も \(\Sigma_0\) の正準自由度に含まれる。完全な隠れ変数を \(\Lambda\) とすると、整合支持が設定に依存するため一般に

\[
\mu_*
\left(
d\Lambda
\mid
x,y,M=\mathsf{Bell}
\right)
\neq
\mu_*
\left(
d\Lambda
\mid
M=\mathsf{Bell}
\right)
\]

となる。Bell の前提違反は測定設定独立性にある \cite{bell1964,chsh1969,wharton2010,wharton_argaman2020,hall2010,leifer_pusey2017,wood_spekkens2015,price_wharton2023,price_wharton2024,argaman2010,hossenfelder_palmer2020}。局所応答、結果の一意性、理想非信号周辺は維持する。Bell の定理を否定しない。

本稿が求めるのは観測条件付き確率 \(P(A,B\mid x,y)\) である。周期内部の設定生成器を外部から差し替え、設定だけを介入変更した場合に同じ分布が保たれるとは示していない。二側境界模型を、前周期記憶が源と設定の共通原因になる前向き模型と同一視もしない。

\hypertarget{ux57faux6e96ux5bc6ux5ea6ux3068ux6709ux9650ux5e45ux306eux8aa4ux5dee}{%
\section{基準密度と有限幅の誤差}\label{ux57faux6e96ux5bc6ux5ea6ux3068ux6709ux9650ux5e45ux306eux8aa4ux5dee}}

基準密度を

\[
q_{AB}^{xy}
=
q^{xy}
\left(
1+\delta_{AB}^{xy}
\right)
\]

とし、

\[
\varepsilon_{\rm base}
=
\max_{A,B,x,y}
\left|
\delta_{AB}^{xy}
\right|
\]

とする。\(\varepsilon_{\rm base}<1\) なら規格化共同分布の理想余弦則からの偏差は \(O(\varepsilon_{\rm base})\) である。

枝作用、境界角、整合判定を含む実際の未規格化質量を \(W_{AB}^{xy}\) とし、理想共通密度を \(c^{xy}K_{AB}^{xy}\) とする。設定対ごとの総偏差を

\[
\varepsilon_{\rm Bell}^{xy}
=
\frac{
\sum_{A,B}
\left|
W_{AB}^{xy}
-
c^{xy}K_{AB}^{xy}
\right|
}{
c^{xy}\mathcal K
}
\]

と定める。\(\varepsilon_{\rm Bell}^{xy}<1\) なら、規格化後の共同分布と理想分布の全変動距離は

\[
d_{\rm TV}^{xy}
\leq
\frac{
\varepsilon_{\rm Bell}^{xy}
}{
1-
\varepsilon_{\rm Bell}^{xy}
}
\]

で抑えられる。位相雑音、戻り損失、枝作用計算誤差、境界角の非一様性、有限比較幅、基準密度非対称性を分けて監査する。

\hypertarget{ux6bd4ux8f03ux5668ux8a18ux9332ux9006ux8a08ux7b97}{%
\section{比較器、記録、逆計算}\label{ux6bd4ux8f03ux5668ux8a18ux9332ux9006ux8a08ux7b97}}

共通未来比較器は、設定と戻りモードから \(K_{AB}^{xy}\) と4区間の境界を可逆に計算する。比較器は過去へ制御信号を送らず、局所記録を変更しない。二側境界問題の終端自由度と整合レジスターを与える。

結果を空の外部記録へ可逆にコピーした後、区間計算、枝作用計算、戻り伝播、局所分析器を逆実行できる。異なる結果を外部記録ごと同じ位相点へ戻すことは Hamiltonian flowの1対1性に反する \cite{landauer1961,bennett1982}。反復には外部記録、不要情報モード、仕事源、弱く結合した環境が必要である。

\hypertarget{ux4e3bux5f35ux306eux7bc4ux56f2-1}{%
\section{主張の範囲}\label{ux4e3bux5f35ux306eux7bc4ux56f2-1}}

本章で得たものは次である。

\begin{enumerate}
\def\labelenumi{\arabic{enumi}.}
\tightlist
\item
  固定左右位相担体と反対称源からの余弦枝作用。
\item
  Bell共同確率を同じ母測度の設定条件付き履歴事象として定義する統一形式。
\item
  任意の基準密度 \(q_{AB}^{xy}\) を含む一般共同確率式。
\item
  共通基準密度下の余弦共同確率、非信号周辺、CHSH値。
\item
  Bell の測定設定独立性が成立しないことの明示。
\end{enumerate}

一方、次は示していない。

\begin{enumerate}
\def\labelenumi{\arabic{enumi}.}
\tightlist
\item
  一般初期集団から反対称源を反復準備すること。
\item
  連続した局所状態から安定な唯一結果を形成すること。
\item
  整合しない周期を事後選別せず零測度にする Hamiltonian 境界値問題。
\item
  4結果の基準密度を共通にする物理的準備。
\item
  目的のBell条件付き部分を持つ不変母測度の具体的な生成。
\item
  外部設定介入に対する統計の安定性。
\item
  一般複合状態、一般測定、一般 Tsirelson 原理。
\end{enumerate}

\hypertarget{ux8aa4ux5deeux53cdux8a3cux6761ux4ef6ux672aux89e3ux6c7aux554fux984c}{%
\chapter{誤差、反証条件、未解決問題}\label{ux8aa4ux5deeux53cdux8a3cux6761ux4ef6ux672aux89e3ux6c7aux554fux984c}}

\begin{statusbox}
位置づけ：現行模型の誤差を、母測度、再帰、相関閉包、正準選択器、作用規格化、有限比較、Bell履歴整合、基準セクター密度、完全周期に分ける。不変性とエルゴード性は仮定であり、具体的な全 Hamiltonian からの導出は主張しない。
\end{statusbox}

\hypertarget{ux8aa4ux5deeux306eux5206ux96e2}{%
\section{誤差の分離}\label{ux8aa4ux5deeux306eux5206ux96e2}}

現行模型の主要誤差と未達条件を次のように分ける。

\begin{longtable}[]{@{}
  >{\raggedright\arraybackslash}p{(\columnwidth - 4\tabcolsep) * \real{0.3333}}
  >{\raggedright\arraybackslash}p{(\columnwidth - 4\tabcolsep) * \real{0.3333}}
  >{\raggedright\arraybackslash}p{(\columnwidth - 4\tabcolsep) * \real{0.3333}}@{}}
\toprule\noalign{}
\begin{minipage}[b]{\linewidth}\raggedright
記号
\end{minipage} & \begin{minipage}[b]{\linewidth}\raggedright
内容
\end{minipage} & \begin{minipage}[b]{\linewidth}\raggedright
中心結論への影響
\end{minipage} \\
\midrule\noalign{}
\endhead
\bottomrule\noalign{}
\endlastfoot
\(\varepsilon_{\rm inv}\) & 周期写像による母測度不変性の破れ & 試行ごとに統計状態を漂移させる \\
\(\varepsilon_{\rm ret}\) & 開始面へ戻らない周期または多重計数 & 母数と完了数をずらし、事後選別を生む \\
\(\varepsilon_{\rm corr}\) & 試行間自己相関と有限混合時間 & 有限試行数の誤差を独立標本評価からずらす \\
\(\varepsilon_{\rm src}\) & 源の相対振幅・相対位相準備誤差 & 初期相関方向、全作用、Bell可視度をずらす \\
\(\varepsilon_{\rm rank}\) & 最大固有値以外の相対重量 & 節の残留強度を変える \\
\(\varepsilon_{\rm close}\) & 非線形項または試行依存 \(h\) による相関閉包誤差 & 交換子発展と階数保存を破る \\
\(\varepsilon_{\rm loc}\) & 空間グラフの不要な遠距離結合 & 局所性と意図した \(h_L\) を破る \\
\(\varepsilon_{\rm disc}\) & 離散演算子と連続演算子の差 & Schrödinger 型連続近似をずらす \\
\(\varepsilon_{\rm sel}\) & 選択器角の条件付き一様性の破れ & 作用区間長と結果頻度の一致を壊す \\
\(\varepsilon_{\rm norm}\) & 全作用と正規化作用分配の相関 & 比の平均を \(C\) の平均比からずらす \\
\(\varepsilon_{\rm cmp}\) & 累積、比較、ラッチの有限幅誤差 & 唯一結果と区間境界をずらす \\
\(\varepsilon_{\rm base}\) & Bell基準セクター密度の非共通性 & 余弦共同分布と非信号周辺をずらす \\
\(\varepsilon_{\rm hist}\) & Bell整合支持条件の不成立 & 不整合周期または事後選別を生む \\
\(\varepsilon_{\rm Bell}\) & 枝作用、境界角、基準密度の総偏差 & 共同分布、非信号性、CHSH値をずらす \\
\(\varepsilon_{\rm cyc}\) & 記録、逆計算、不要情報移送、再準備の周期誤差 & 反復可能性と母測度の再現性を破る \\
\end{longtable}

これらは独立とは限らない。それでも、どの段階で生じた誤差かを記録し、1つの総誤差だけで原因を隠さない。

\hypertarget{ux6bcdux6e2cux5ea6ux518dux5e30ux9577ux671fux983bux5ea6}{%
\section{母測度、再帰、長期頻度}\label{ux6bcdux6e2cux5ea6ux518dux5e30ux9577ux671fux983bux5ea6}}

第1章では

\[
\mathcal R_*\mu_*
=
\mu_*,
\qquad
\mu_*(\tau<\infty)=1
\]

と、\(\mu_*\) に関するエルゴード性を仮定した。これらは相関行列または選択器の有限次元代数から従わない。

有限データで不変性を検査するには、試行観測量 \(f_j\) の集合に対し

\[
\varepsilon_{\rm inv}
=
\max_j
\left|
\mathbb E_{\mu_*}
\left[
f_j\circ\mathcal R
\right]
-
\mathbb E_{\mu_*}
\left[
f_j
\right]
\right|
\]

を測る。有限個の観測量で零でも、測度全体の不変性を証明したことにはならない。

エルゴード性は長期平均の極限を与えるが、有限試行数での収束率を与えない。結果指示関数、全作用、選択器角、設定、源状態の自己相関を測り、有効標本数を評価する必要がある。混合が遅くてもエルゴード的であり得るため、\(\varepsilon_{\rm corr}\) をエルゴード性の真偽と混同しない。

\hypertarget{ux6b63ux6e96ux9078ux629eux5668ux306eux8aa4ux5dee}{%
\section{正準選択器の誤差}\label{ux6b63ux6e96ux9078ux629eux5668ux306eux8aa4ux5dee}}

第4章の一般式は

\[
P_k
=
\mathbb E
\left[
\frac{I_k}{I_{\rm ph}}
\right]
\]

である。選択器角の周辺一様性だけでは足りない。作用分配、測定基底、調製条件の下での条件付き分布を検査する。

全作用変動の補正は

\[
\varepsilon_{{\rm norm},k}
=
\left|
\frac{
\operatorname{Cov}
\left(
I_{\rm ph},
I_k/I_{\rm ph}
\right)
}{
\mathbb E[I_{\rm ph}]
}
\right|
\]

である。全作用の分散が小さくても、正規化分配との相関が強ければ無視できない。

比較器の半幅を \(w\) とすると、境界近傍の母測度

\[
\varepsilon_{\rm cmp}(w)
=
\mu_*
\left(
\min_k
\left|u-S_k\right|
\leq w
\middle|
M=\mathsf{basis}
\right)
\]

を測る。\(w\) を減らした収束、ラッチの二重作動率、無結果率、結果別失敗率を別に記録する。

\hypertarget{ux76f8ux95a2ux9589ux5305ux3068ux7a7aux9593ux7d50ux5408}{%
\section{相関閉包と空間結合}\label{ux76f8ux95a2ux9589ux5305ux3068ux7a7aux9593ux7d50ux5408}}

閉包残差 \(D_C\) を

\[
i\mathcal J_0\dot C
-
\left[h,C\right]
=
D_C
\]

で定める。有限観測時間 \([0,T]\) の積分誤差を

\[
\varepsilon_{\rm close}(T)
=
\frac1{\mathcal J_0}
\int_0^T
\left\|D_C(t)\right\|_{\rm op}
\,dt
\]

とする。同じ条件付き集団の中で \(h\) が試行ごとに異なる場合も \(D_C\) に含める。4次以上の項がある場合は、高次モーメント、振幅範囲、観測時間を変えて元の有限試行を直接検証する。

\(h_L\) の物理的実装では、不要な遠距離結合、\(Q\) と \(P\) の結合行列の差、境界条件、格子幅収束を監査する。\(\mathcal J_0^2/(2m)\) は現模型の設計係数であり、調整後に Schrödinger 型方程式と一致することは、その係数の創発的導出ではない。

\hypertarget{ux6e90ux6e96ux5099ux968eux6570ux6b20ux9665ux5e72ux6e09}{%
\section{源準備、階数欠陥、干渉}\label{ux6e90ux6e96ux5099ux968eux6570ux6b20ux9665ux5e72ux6e09}}

理想出力を \(e^{i\beta^\omega}\chi_0\) とし、実出力を

\[
b^\omega
=
e^{i\beta^\omega}\chi_0
+
\delta b^\omega
\]

と書く。共通位相方向を除いた \(\delta b^\omega\) の2次平均を源準備誤差とする。絶対位相が一様でも \(C\) は階数1になり得るため、\(\mathbb E[b]\) の小ささを準備失敗と判定しない。

階数欠陥については、第3章の節上界

\[
p_{\rm node}
\leq
\varepsilon_{\rm rank}
\]

を使う。再結合器誤差と位相雑音を含む場合は、節の絶対強度と全変動距離で評価する。閉鎖線形発展では \(\varepsilon_{\rm rank}\) は保存される。

\hypertarget{bell-ux5c65ux6b74ux8aa4ux5dee}{%
\section{Bell 履歴誤差}\label{bell-ux5c65ux6b74ux8aa4ux5dee}}

Bell側の一般式は

\[
P(A,B\mid x,y,G)
=
\frac{
q_{AB}^{xy}K_{AB}^{xy}
}{
\sum_{A',B'}
q_{A'B'}^{xy}K_{A'B'}^{xy}
}
\]

である。従って、余弦枝作用 \(K_{AB}^{xy}\) が正しくても、\(q_{AB}^{xy}\) が不均等なら理想共同分布は得られない。基準密度の結果依存性を、枝作用計算誤差へ吸収して隠さない。

整合支持の失敗率を

\[
\varepsilon_{\rm hist}
=
\mu_*
\left(
G^c
\mid
M=\mathsf{Bell}
\right)
\]

とする。本模型の理想仮定では \(\varepsilon_{\rm hist}=0\) である。実験で \(G^c\) を捨てて零に見せることは許されない。開始数、設定生成数、左右局所記録数、比較完了数、外部記録数、次周期再帰数を設定別・結果別に監査する。

\hypertarget{ux7c92ux5b50ux4f4dux7f6eux3068ux306eux672aux63a5ux7d9a}{%
\section{粒子位置との未接続}\label{ux7c92ux5b50ux4f4dux7f6eux3068ux306eux672aux63a5ux7d9a}}

第4章の結果 \(E_k^U\) は位相担体モードの選択事象である。粒子位置の主張には、少なくとも次が必要である。

\begin{enumerate}
\def\labelenumi{\arabic{enumi}.}
\tightlist
\item
  粒子座標 \(X\) と局所検出器の明示 Hamiltonian。
\item
  選択された位相担体モードと局所粒子反応の同期。
\item
  1試行1粒子結果を保つ増幅と記録。
\item
  途中の粒子流束または連続軌道。
\item
  2重スリット検出面での局所反応。
\end{enumerate}

これらがない限り、\(P_i=C_{ii}/\operatorname{tr}C\) を粒子の位置 Born 則が完全に導出されたとは評価しない。

\hypertarget{bell-ux306eux524dux63d0ux76e3ux67fb}{%
\section{Bell の前提監査}\label{bell-ux306eux524dux63d0ux76e3ux67fb}}

\begin{longtable}[]{@{}
  >{\raggedright\arraybackslash}p{(\columnwidth - 2\tabcolsep) * \real{0.5000}}
  >{\raggedright\arraybackslash}p{(\columnwidth - 2\tabcolsep) * \real{0.5000}}@{}}
\toprule\noalign{}
\begin{minipage}[b]{\linewidth}\raggedright
前提
\end{minipage} & \begin{minipage}[b]{\linewidth}\raggedright
現行模型での扱い
\end{minipage} \\
\midrule\noalign{}
\endhead
\bottomrule\noalign{}
\endlastfoot
局所応答 & \(A=A(x,\Lambda_A)\)、\(B=B(y,\Lambda_B)\) を維持する \\
結果の一意性 & 理想局所結果セクターとして仮定する。一般形成機構は未完成 \\
測定設定独立性 & 成立しない。\(\mu_*(d\Lambda\mid x,y)\) は設定に依存し得る \\
事後選別なし & \(\mu_*(G\mid M=\mathsf{Bell})=1\) を要求する。物理的生成は未完成 \\
非信号性 & 共通基準密度と理想余弦枝作用の下で成立する \\
試行測度 & 相関、有限基底、Bellで共通の \(\mu_*\) を用いる \\
\end{longtable}

Bell の定理は正しい。本模型がCHSH不等式を破る条件付き相関を持てるのは、Bell の測定設定独立性を満たさないからである。

\hypertarget{ux53cdux8a3cux6761ux4ef6}{%
\section{反証条件}\label{ux53cdux8a3cux6761ux4ef6}}

次のいずれかが示されれば、対応する現行主張は成立しない。

\begin{enumerate}
\def\labelenumi{\arabic{enumi}.}
\tightlist
\item
  同じ条件付き集団の位相担体が共通2次 Hamiltonian に従わず、閉包残差が観測時間で小さくならない。
\item
  周期写像が目的の母測度を保存しない、または結果別に戻らない周期が残る。
\item
  選択器角が作用分配の下で条件付き一様にならず、偏りが装置極限で消えない。
\item
  固定作用セクターを準備できず、共分散補正が観測精度より大きい。
\item
  累積比較器の二重結果率または無結果率が比較幅を狭めても消えない。
\item
  単粒子側で測定基底を変えると調製相関 \(C_U\) が変わり、同じ状態の基底測定と解釈できない。
\item
  Bell基準密度 \(q_{AB}^{xy}\) の非対称性が残り、余弦共同分布または非信号周辺が崩れる。
\item
  Bell不整合周期を捨てなければ目標分布が出ない。
\item
  非信号性が結果依存損失または設定依存の検出効率に依存する。
\item
  不要情報を含む拡大全系で、再準備写像を1対1の Hamiltonian 流として実装できない。
\end{enumerate}

\hypertarget{ux65e7ux6a21ux578bux3092ux5916ux3057ux305fux3053ux3068ux306bux3088ux308bux5909ux66f4}{%
\section{旧模型を外したことによる変更}\label{ux65e7ux6a21ux578bux3092ux5916ux3057ux305fux3053ux3068ux306bux3088ux308bux5909ux66f4}}

本改訂では、位置結果の主機構から2モード作用殻、共通入口流束、局所作用の排他的全入口転送を外した。これにより、階数1だけでなく高階数相関行列と任意有限基底を扱える。一方、正準角の条件付き一様性と累積比較器が新しい成立条件になった。

Bell側では、2モード比較殻の容量を確率機構から外し、枝作用に比例する区間と二側履歴整合へ置き換えた。しかし、4結果で共通な境界密度の課題は消えていない。新しい一般式では \(q_{AB}^{xy}\) として明示される。

旧2成分誘導場を外したため、循環量子化、粒子・場同期、位相接続による粒子誘導は現行結果に含まれない。Wallstrom問題と粒子の連続運動は未解決である。

\hypertarget{ux6700ux91cdux8981ux306eux672aux89e3ux6c7aux554fux984c}{%
\section{最重要の未解決問題}\label{ux6700ux91cdux8981ux306eux672aux89e3ux6c7aux554fux984c}}

\begin{enumerate}
\def\labelenumi{\arabic{enumi}.}
\tightlist
\item
  \(h_L\) の局所結合構造と係数を、量子力学的演算子の転記でなく、より原始的な有限古典模型から導く。
\item
  目的の \(\mu_*\) を持つ具体的な周期写像を全 Hamiltonian から構成し、不変性、再帰性、仮定したエルゴード性を検証する。
\item
  選択器角の条件付き一様性を、測定結合と再準備を含む周期で維持する。
\item
  累積作用レジスター、有限幅比較器、唯一結果ラッチ、外部記録、逆計算を1つの滑らかな有限 Hamiltonian へ統合する。
\item
  位相担体の有限基底結果を粒子の局所検出と連続運動へ接続する。
\item
  非線形項を含む有限時間の相関閉包誤差を制御する。
\item
  空間グラフの閉路巻数、節、位相すり抜け、連続補間を統合し、Wallstrom問題へ回答する。
\item
  井戸型・調和型ポテンシャル、トンネル効果、Rabi振動を定量検証する。
\item
  反対称 Bell源を反復可能な有限 Hamiltonian装置で準備し、相対位相を保持する。
\item
  Bell整合支持条件と共通基準セクター密度を、事後選別なしの二側境界周期として実現する。
\item
  非対称準備と有限誤差を含む一般的な非信号条件を示す。
\item
  平面内2出力を超える一般複合状態と局所測定へ拡張する。
\end{enumerate}

\appendix

\hypertarget{ux6b63ux6e96ux4f4dux76f8ux62c5ux4f53ux76f8ux95a2ux884cux5217ux968eux65701ux56e0ux5b50ux306eux8a3cux660e}{%
\chapter{正準位相担体、相関行列、階数1因子の証明}\label{ux6b63ux6e96ux4f4dux76f8ux62c5ux4f53ux76f8ux95a2ux884cux5217ux968eux65701ux56e0ux5b50ux306eux8a3cux660e}}

\begin{statusbox}
位置づけ：本文第2章と第3章で用いる複素正準表示、相関行列発展、相関スペクトル保存、階数1の同値条件、共通源準備、階数1因子、近似階数1の節上界を有限次元で証明する。
\end{statusbox}

\hypertarget{ux8907ux7d20ux6b63ux6e96ux5ea7ux6a19}{%
\section{複素正準座標}\label{ux8907ux7d20ux6b63ux6e96ux5ea7ux6a19}}

実正準対 \((Q_i,P_i)\) と固定作用尺度 \(\mathcal J_0>0\) に対し、

\[
b_i
=
\frac{Q_i+iP_i}{\sqrt{2\mathcal J_0}},
\qquad
b_i^*
=
\frac{Q_i-iP_i}{\sqrt{2\mathcal J_0}}
\]

とする。直接計算により

\[
\left\{b_i,b_j\right\}
=
0,
\qquad
\left\{b_i^*,b_j^*\right\}
=
0,
\]

\[
\left\{b_i,b_j^*\right\}
=
-
\frac{i}{\mathcal J_0}
\delta_{ij}
\]

を得る。

実関数 \(H(b,b^*)\) に対する Hamilton 方程式は

\[
\dot b_i
=
-
\frac{i}{\mathcal J_0}
\frac{\partial H}{\partial b_i^*}
\]

である。\(H=b^\dagger hb\) なら \(\partial H/\partial b^*=hb\) なので、

\[
i\mathcal J_0\dot b
=
hb
\]

となる。

\hypertarget{ux5b9f2ux6b21-hamiltonian-ux306eux8868ux793a}{%
\section{実2次 Hamiltonian の表示}\label{ux5b9f2ux6b21-hamiltonian-ux306eux8868ux793a}}

\(h=A+iB\) と分ける。Hermitian 条件から

\[
A^{\mathsf T}=A,
\qquad
B^{\mathsf T}=-B
\]

である。\(Q=(Q_1,\ldots,Q_L)^{\mathsf T}\)、\(P=(P_1,\ldots,P_L)^{\mathsf T}\) とすると、

\[
b^\dagger hb
=
\frac1{2\mathcal J_0}
\left[
Q^{\mathsf T}AQ
+
P^{\mathsf T}AP
-
2Q^{\mathsf T}BP
\right]
\]

である。右辺は実正準変数の通常の2次 Hamiltonian である。\(h\) が実対称なら \(B=0\) で、\(Q\) と \(P\) に同じ2次形式が作用する。

\hypertarget{ux5171ux901aux56deux8ee2ux3068ux5168ux4f4dux76f8ux4f5cux7528}{%
\section{共通回転と全位相作用}\label{ux5171ux901aux56deux8ee2ux3068ux5168ux4f4dux76f8ux4f5cux7528}}

全位相作用を

\[
I_{\rm ph}
=
\frac12
\left(
Q^{\mathsf T}Q
+
P^{\mathsf T}P
\right)
=
\mathcal J_0b^\dagger b
\]

とする。Poisson 括弧は

\[
\left\{b_i,I_{\rm ph}\right\}
=
-ib_i
\]

なので、\(I_{\rm ph}\) は共通位相回転を生成する。

\(H=b^\dagger hb\) に対し、

\[
\left\{I_{\rm ph},H\right\}
=
0
\]

である。従って、時間依存 \(h(t)\) であっても共通回転対称性が保たれる限り \(I_{\rm ph}\) は保存される。

\hypertarget{ux76f8ux95a2ux884cux5217ux65b9ux7a0bux5f0f}{%
\section{相関行列方程式}\label{ux76f8ux95a2ux884cux5217ux65b9ux7a0bux5f0f}}

各試行で

\[
i\mathcal J_0\dot b^\omega
=
h(t)b^\omega
\]

とする。積を微分すると、

\[
\begin{aligned}
i\mathcal J_0
\frac{d}{dt}
\left(
b^\omega
\left(b^\omega\right)^\dagger
\right)
={}&
h
b^\omega
\left(b^\omega\right)^\dagger
\\
&-
b^\omega
\left(b^\omega\right)^\dagger
h
\end{aligned}
\]

となる。母測度 \(\mu_*\) を調製条件 \(\mathcal P\) とプログラム \(M\) で条件付ける。同じ条件付き集団の全試行が共通の \(h(t)\) に従い、条件が観測窓で固定され、微分と条件付き平均を交換できるなら、

\[
i\mathcal J_0\dot C
=
\left[h,C\right]
\]

を得る。ここで

\[
C(t)
=
\mathbb E_{\mu_*}
\left[
b_t b_t^\dagger
\mid
\mathcal P,M
\right]
\]

である。母測度の不変性またはエルゴード性は、この有限時間微分の代数には不要である。これらは同じ条件付き集団を周期間で再現し、長期頻度と母測度確率を結ぶために使う。

\hypertarget{ux76f8ux95a2ux30b9ux30daux30afux30c8ux30ebux306eux4fddux5b58}{%
\section{相関スペクトルの保存}\label{ux76f8ux95a2ux30b9ux30daux30afux30c8ux30ebux306eux4fddux5b58}}

時間発展作用素 \(U\) は

\[
i\mathcal J_0\dot U
=
hU,
\qquad
U(t_0)=I
\]

を満たす。\(h=h^\dagger\) なら

\[
\frac{d}{dt}
\left(U^\dagger U\right)
=
0
\]

なので \(U\) はユニタリである。相関行列は

\[
C(t)
=
U(t,t_0)
C(t_0)
U(t,t_0)^\dagger
\]

となる。

従って、任意の正整数 \(n\) に対し

\[
\operatorname{tr}C(t)^n
=
\operatorname{tr}C(t_0)^n
\]

である。有限次元では固有値の全対称多項式が保存されるため、固有値の多重集合、階数、純度が保存される。

\hypertarget{ux968eux65701ux76f8ux95a2ux306aux3089ux5171ux901aux5c04ux5f71ux65b9ux5411}{%
\section{階数1相関なら共通射影方向}\label{ux968eux65701ux76f8ux95a2ux306aux3089ux5171ux901aux5c04ux5f71ux65b9ux5411}}

\(C=\Lambda\chi\chi^\dagger\)、\(\Lambda>0\)、\(\chi^\dagger\chi=1\) とする。\(\chi\) と直交する任意の \(v\) に対し、

\[
0
=
v^\dagger Cv
=
\mathbb E
\left[
\left|v^\dagger b\right|^2
\right]
\]

である。非負確率変数の平均が零なので、

\[
v^\dagger b^\omega
=
0
\]

がほとんど確実に成立する。

有限次元直交補空間の正規直交基底 \(v_2,\ldots,v_L\) を取ると、全てについて同時に上式が成立する確率1の集合を取れる。従って

\[
b^\omega
=
c^\omega\chi
\]

がほとんど確実に成立する。ここで \(c^\omega=\chi^\dagger b^\omega\) である。さらに

\[
C
=
\mathbb E
\left|c^\omega\right|^2
\chi\chi^\dagger
\]

なので \(\Lambda=\mathbb E|c^\omega|^2\) である。

逆に \(b^\omega=c^\omega\chi\) がほとんど確実なら、上式から \(C\) は階数1である。これで本文第2.6節の同値条件が従う。

\hypertarget{ux5171ux901aux6e90ux6e96ux5099}{%
\section{共通源準備}\label{ux5171ux901aux6e90ux6e96ux5099}}

入力が

\[
b_{\rm in}^\omega
=
c^\omega e_0
\]

であり、準備回路 \(U_{\rm prep}\) が \(U_{\rm prep}e_0=\chi_0\) を満たすとする。線形性により

\[
b_{\rm out}^\omega
=
c^\omega\chi_0
\]

である。従って

\[
C_{\rm out}
=
\mathbb E
\left|c^\omega\right|^2
\chi_0\chi_0^\dagger
\]

となる。

\(c^\omega=e^{i\beta^\omega}\) で \(\beta^\omega\) が一様なら、

\[
\mathbb E
\left[b_{\rm out}^\omega\right]
=
0
\]

であっても、

\[
C_{\rm out}
=
\chi_0\chi_0^\dagger
\]

である。1次平均でなく2次相関を状態量に選ぶ理由がここにある。

\hypertarget{ux968eux65701ux56e0ux5b50ux306eux767aux5c55}{%
\section{階数1因子の発展}\label{ux968eux65701ux56e0ux5b50ux306eux767aux5c55}}

\(C=\Lambda\chi\chi^\dagger\)、\(\chi^\dagger\chi=1\) とし、\(\Lambda\) は一定とする。交換子方程式へ代入すると、

\[
i\mathcal J_0
\left(
\dot\chi\chi^\dagger
+
\chi\dot\chi^\dagger
\right)
=
h\chi\chi^\dagger
-
\chi\chi^\dagger h
\]

である。

左から \(I-\chi\chi^\dagger\)、右から \(\chi\) を掛けると、

\[
\left(
I-\chi\chi^\dagger
\right)
\left(
i\mathcal J_0\dot\chi-h\chi
\right)
=
0
\]

を得る。従って

\[
i\mathcal J_0\dot\chi
-
h\chi
=
\lambda(t)\chi
\]

である。規格化を微分し、\(h\) が Hermitian であることを用いると \(\lambda(t)\) は実数となる。

位相変換

\[
\widetilde\chi(t)
=
\exp
\left[
\frac{i}{\mathcal J_0}
\int_{t_0}^t
\lambda(s)\,ds
\right]
\chi(t)
\]

により、

\[
i\mathcal J_0
\dot{\widetilde\chi}
=
h\widetilde\chi
\]

となる。

\hypertarget{ux8fd1ux4f3cux968eux65701ux3068ux7bc0ux4e0aux754c}{%
\section{近似階数1と節上界}\label{ux8fd1ux4f3cux968eux65701ux3068ux7bc0ux4e0aux754c}}

\(C\) の最大固有値を \(\lambda_1\)、単位主固有ベクトルを \(\chi\) とする。スペクトル分解から

\[
C
=
\lambda_1\chi\chi^\dagger
+
E,
\qquad
E\geq0,
\qquad
E\chi=0
\]

である。\(\operatorname{tr}E=\operatorname{tr}C-\lambda_1\) なので、

\[
\varepsilon_{\rm rank}
=
\frac{\operatorname{tr}E}{\operatorname{tr}C}
\]

となる。

任意のユニタリ \(U\) と出力基底ベクトル \(e_k\) に対し、\((U\chi)_k=0\) なら

\[
\begin{aligned}
p_k
&=
\frac{
e_k^\dagger UCU^\dagger e_k
}{
\operatorname{tr}C
}
\\
&=
\frac{
e_k^\dagger UEU^\dagger e_k
}{
\operatorname{tr}C
}
\\
&\leq
\frac{\left\|E\right\|_{\rm op}}{\operatorname{tr}C}
\leq
\frac{\operatorname{tr}E}{\operatorname{tr}C}
=
\varepsilon_{\rm rank}.
\end{aligned}
\]

正半定値性が最後の不等式に必要である。

\hypertarget{ux9589ux5305ux6b8bux5deeux306eux4e0aux754c}{%
\section{閉包残差の上界}\label{ux9589ux5305ux6b8bux5deeux306eux4e0aux754c}}

残差付き試行方程式

\[
i\mathcal J_0\dot b
=
hb+r
\]

から

\[
D_C
=
\mathbb E
\left[
rb^\dagger-br^\dagger
\right]
\]

を得る。任意の単位ベクトル \(u,v\) に対し、

\[
\left|
u^\dagger
\mathbb E
\left[rb^\dagger\right]
v
\right|
\leq
\left(
\mathbb E\left\|r\right\|^2
\right)^{1/2}
\left(
\mathbb E\left\|b\right\|^2
\right)^{1/2}
\]

なので、2項を合わせて

\[
\left\|D_C\right\|_{\rm op}
\leq
2
\left(
\mathbb E\left\|r\right\|^2
\right)^{1/2}
\left(
\mathbb E\left\|b\right\|^2
\right)^{1/2}
\]

となる。これは残差の起源を導く式ではなく、試行残差から相関閉包誤差への伝播上界である。

\hypertarget{ux6bcdux6e2cux5ea6ux4f5cux7528ux9078ux629eux5668ux5c65ux6b74ux533aux9593ux306eux8a3cux660e}{%
\chapter{母測度、作用選択器、履歴区間の証明}\label{ux6bcdux6e2cux5ea6ux4f5cux7528ux9078ux629eux5668ux5c65ux6b74ux533aux9593ux306eux8a3cux660e}}

\begin{statusbox}
位置づけ：本文第1章、第4章、第5章で用いる母測度の条件付け、エルゴード頻度、作用区間選択、固定作用公式、共分散補正、任意有限基底、Bell履歴区間と基準セクター密度を補足する。不変母測度の力学的生成は仮定である。
\end{statusbox}

\hypertarget{ux6761ux4ef6ux4ed8ux304dux6bcdux6e2cux5ea6}{%
\section{条件付き母測度}\label{ux6761ux4ef6ux4ed8ux304dux6bcdux6e2cux5ea6}}

試行開始面を \((\Sigma_0,\mathcal B)\)、母測度を \(\mu_*\) とする。調製条件、プログラム、基底、Bell設定は全て \(\mathcal B\) 可測な開始点の関数またはレジスター領域として扱う。

条件事象 \(D\in\mathcal B\)、\(\mu_*(D)>0\) に対し、

\[
\mu_*^D(E)
=
\frac{
\mu_*(E\cap D)
}{
\mu_*(D)
}
\]

とする。本文の \(\mathbb E_{\mu_*}[\cdot\mid\mathcal P,M,U]\) は、連続レジスターでは正則条件付き分布、離散レジスターでは上の条件事象による期待値を表す。

相関行列は

\[
C_D(t)
=
\int_{\Sigma_0}
b_t(z)b_t(z)^\dagger
\,d\mu_*^D(z)
\]

である。有限基底結果とBell結果も同じ \(\mu_*^D\) の事象確率であるため、3つに独立な基礎測度を置かない。

\hypertarget{ux4e0dux5909ux6027ux3068ux30a8ux30ebux30b4ux30fcux30c9ux983bux5ea6}{%
\section{不変性とエルゴード頻度}\label{ux4e0dux5909ux6027ux3068ux30a8ux30ebux30b4ux30fcux30c9ux983bux5ea6}}

周期写像 \(\mathcal R\) が

\[
\mathcal R_*\mu_*
=
\mu_*
\]

を満たし、\(\mu_*\) に関してエルゴード的であると仮定する。Birkhoffのエルゴード定理により、\(f\in L^1(\mu_*)\) に対して

\[
\frac1N
\sum_{n=0}^{N-1}
f\left(
\mathcal R^nz
\right)
\longrightarrow
\int f\,d\mu_*
\]

が \(\mu_*\) に関してほとんど確実に成立する。

事象 \(E\) の指示関数を取れば

\[
\frac1N
\sum_{n=0}^{N-1}
\mathbf 1_E
\left(
\mathcal R^nz
\right)
\longrightarrow
\mu_*(E)
\]

となる。条件事象 \(D\) が正の測度を持ち、その出現回数が無限に増えるなら、比エルゴード定理または分子と分母への上式の適用により

\[
\frac{
\sum_{n=0}^{N-1}
\mathbf 1_{E\cap D}
\left(
\mathcal R^nz
\right)
}{
\sum_{n=0}^{N-1}
\mathbf 1_D
\left(
\mathcal R^nz
\right)
}
\longrightarrow
\mu_*(E\mid D)
\]

を得る。これは設定対ごとのBell相対頻度に用いる。エルゴード性は本稿の仮定であり、以下の有限次元積分からは導かれない。

\hypertarget{ux4f5cux7528ux533aux9593ux9078ux629e}{%
\section{作用区間選択}\label{ux4f5cux7528ux533aux9593ux9078ux629e}}

\(I_k\geq0\)、\(I_{\rm ph}=\sum_kI_k>0\) とする。\(\vartheta\) が \((I_1,\ldots,I_L)\) の下で \([0,2\pi)\) 上に条件付き一様であるとする。

\[
u
=
\frac{\vartheta}{2\pi}
I_{\rm ph},
\qquad
S_k
=
\sum_{j=1}^kI_j
\]

とし、\(E_k=\{S_{k-1}\leq u<S_k\}\) とする。

\begin{lemma}[作用区間の条件付き確率]

$$
P(E_k\mid I_1,\ldots,I_L)
=
\frac{I_k}{I_{\rm ph}}
$$

が成立する。
\end{lemma}

\begin{proof}
$u$ の条件付き密度は $1/I_{\rm ph}$ である。区間 $[S_{k-1},S_k)$ の長さが $I_k$ なので、積分は $I_k/I_{\rm ph}$ となる。
\end{proof}

\(I_k=0\) の区間は空であり、正の確率を持たない。境界 \(u=S_k\) の集合は、条件付き一様分布では零測度である。半開区間の選び方は確率を変えない。

\hypertarget{ux6b63ux6e96ux57faux5e95ux6df7ux5408ux3068ux56faux5b9aux4f5cux7528ux516cux5f0f}{%
\section{正準基底混合と固定作用公式}\label{ux6b63ux6e96ux57faux5e95ux6df7ux5408ux3068ux56faux5b9aux4f5cux7528ux516cux5f0f}}

\(b\in\mathbb C^L\)、\(U^\dagger U=I\) とし、

\[
I_k
=
\mathcal J_0
\left|
\left(Ub\right)_k
\right|^2,
\qquad
I_{\rm ph}
=
\mathcal J_0b^\dagger b
\]

とする。ユニタリ性により \(\sum_kI_k=I_{\rm ph}\) である。

条件付き相関行列を

\[
C_U
=
\mathbb E
\left[
bb^\dagger
\mid
U,\mathcal P,M=\mathsf{basis}
\right]
\]

とする。\(I_{\rm ph}=I_0\) がほとんど確実なら、

\[
\begin{aligned}
P_k
&=
\mathbb E
\left[
\frac{I_k}{I_0}
\right]
\\
&=
\frac{
\mathcal J_0
\mathbb E
\left|
\left(Ub\right)_k
\right|^2
}{
I_0
}
\\
&=
\frac{
\left(
UC_UU^\dagger
\right)_{kk}
}{
\operatorname{tr}C_U
}.
\end{aligned}
\]

この導出は \(C_U\) の階数を使わない。

\hypertarget{ux5168ux4f5cux7528ux5909ux52d5ux306eux5171ux5206ux6563ux6052ux7b49ux5f0f}{%
\section{全作用変動の共分散恒等式}\label{ux5168ux4f5cux7528ux5909ux52d5ux306eux5171ux5206ux6563ux6052ux7b49ux5f0f}}

\[
r_k
=
\frac{I_k}{I_{\rm ph}}
\]

とする。\(I_k=I_{\rm ph}r_k\) なので、

\[
\mathbb E[I_k]
=
\mathbb E[I_{\rm ph}]
\mathbb E[r_k]
+
\operatorname{Cov}
\left(
I_{\rm ph},r_k
\right)
\]

である。従って、

\[
\mathbb E[r_k]
-
\frac{
\mathbb E[I_k]
}{
\mathbb E[I_{\rm ph}]
}
=
-
\frac{
\operatorname{Cov}
\left(
I_{\rm ph},r_k
\right)
}{
\mathbb E[I_{\rm ph}]
}
\]

となる。固定作用は共分散を零にする十分条件である。独立性までは不要で、無相関でもよい。

\hypertarget{ux6e2cux5b9aux57faux5e95ux306bux4f9dux5b58ux3057ux306aux3044ux8abfux88fd}{%
\section{測定基底に依存しない調製}\label{ux6e2cux5b9aux57faux5e95ux306bux4f9dux5b58ux3057ux306aux3044ux8abfux88fd}}

異なる \(U\) に対する条件付き相関を \(C_U\) とする。\(C_U=C\) が全ての許容基底で成立するとき、固定作用公式は

\[
P_k(U)
=
\frac{
\left(
UCU^\dagger
\right)_{kk}
}{
\operatorname{tr}C
}
\]

となる。

\(C_U\neq C\) なら同じ代数は

\[
P_k(U)
=
\frac{
\left(
UC_UU^\dagger
\right)_{kk}
}{
\operatorname{tr}C_U
}
\]

を与えるだけである。これは基底依存の調製集団に対する式であり、1つの状態 \(C\) の測定基底を変えた結果とは呼べない。

\hypertarget{ux6761ux4ef6ux4ed8ux304dux4e00ux69d8ux6027ux304bux3089ux306eux305aux308c}{%
\section{条件付き一様性からのずれ}\label{ux6761ux4ef6ux4ed8ux304dux4e00ux69d8ux6027ux304bux3089ux306eux305aux308c}}

理想条件付き角分布を \(m(d\vartheta)=d\vartheta/(2\pi)\)、実分布を \(\rho_{b,U}(d\vartheta)\) とする。任意の区間事象 \(E_k\) に対して、全変動距離の定義から

\[
\left|
\rho_{b,U}(E_k)
-
m(E_k)
\right|
\leq
d_{\rm TV}
\left(
\rho_{b,U},m
\right)
\]

である。母測度で平均すれば、

\[
\left|
P_k^{\rm real}
-
P_k^{\rm ideal}
\right|
\leq
\mathbb E
\left[
d_{\rm TV}
\left(
\rho_{b,U},m
\right)
\right]
\]

を得る。周辺角分布の全変動距離ではなく、\((b,U,\mathcal P)\) で条件付けた距離を使う必要がある。

\hypertarget{ux6709ux9650ux5e45ux6bd4ux8f03ux5668}{%
\section{有限幅比較器}\label{ux6709ux9650ux5e45ux6bd4ux8f03ux5668}}

理想境界集合を

\[
\mathcal D
=
\bigcup_{k=1}^{L-1}
\{u=S_k\}
\]

とする。有限幅比較器が理想結果と異なり得る領域を

\[
\mathcal D_w
=
\left\{
\min_k
\left|u-S_k\right|
\leq w
\right\}
\]

へ限定できるなら、結合不等式により理想分布と実分布の全変動距離は

\[
d_{\rm TV}
\leq
\mu_*(\mathcal D_w)
\]

で抑えられる。

条件付き \(u\) 密度が \(M_u\) 以下なら、粗い上界として

\[
\mu_*(\mathcal D_w)
\leq
2(L-1)M_uw
\]

を得る。ただし、境界が接近または重複する場合は過大評価である。有限装置では二重ラッチ、無結果、比較順序依存も別に測る。

\hypertarget{bellux679dux533aux9593}{%
\section{Bell枝区間}\label{bellux679dux533aux9593}}

固定設定 \((x,y)\) で、4つの非負枝作用 \(K_{AB}^{xy}\) が

\[
\sum_{A,B}
K_{AB}^{xy}
=
\mathcal K
\]

を満たすとする。\(u_{\rm B}\) を \([0,\mathcal K)\) 上の一様変数とし、区間 \(\mathcal I_{AB}^{xy}\) の長さを \(K_{AB}^{xy}\) とする。

Bell側では \((A,B)\) が先に局所記録されるため、区間は結果生成に使わない。候補完結履歴の結果セクターを固定した後、

\[
G_{AB}^{xy}
=
\left\{
u_{\rm B}
\in
\mathcal I_{AB}^{xy}
\right\}
\]

を整合条件とする。

\hypertarget{ux57faux6e96ux30bbux30afux30bfux30fcux5bc6ux5ea6ux3092ux542bux3080-bellux516cux5f0f}{%
\section{基準セクター密度を含む Bell公式}\label{ux57faux6e96ux30bbux30afux30bfux30fcux5bc6ux5ea6ux3092ux542bux3080-bellux516cux5f0f}}

選択器座標を除く候補履歴変数を \(\zeta\) とする。固定設定と結果セクターで、候補 Liouville 基準要素が

\[
d\nu_{AB}^{xy}
=
q_{AB}^{xy}
d\bar\nu_{AB}^{xy}(\zeta)
\frac{du_{\rm B}}{\mathcal K},
\qquad
\int
d\bar\nu_{AB}^{xy}
=
1
\]

と分解できるとする。\(q_{AB}^{xy}\) は選択器以外の密度と体積を全て含む。

整合事象の基準質量は

\[
\nu_{AB}^{xy}
\left(
G_{AB}^{xy}
\right)
=
q_{AB}^{xy}
\frac{K_{AB}^{xy}}{\mathcal K}
\]

である。全4セクターを合わせて規格化すると、

\[
P(A,B\mid x,y,G)
=
\frac{
q_{AB}^{xy}K_{AB}^{xy}
}{
\sum_{A',B'}
q_{A'B'}^{xy}K_{A'B'}^{xy}
}
\]

を得る。

\begin{corollary}[共通基準密度]
$q_{AB}^{xy}=q^{xy}$ が4結果で共通なら、

$$
P(A,B\mid x,y,G)
=
\frac{K_{AB}^{xy}}{\mathcal K}
$$

である。
\end{corollary}

この系は共通性を導くものではない。\(q_{AB}^{xy}\) が不均等なら、一様なBell境界角だけでは余弦共同確率を得られない。

\hypertarget{ux5168ux8a66ux884cux76e3ux67fbux3068ux518dux521dux671fux5316}{%
\section{全試行監査と再初期化}\label{ux5168ux8a66ux884cux76e3ux67fbux3068ux518dux521dux671fux5316}}

母測度確率を実験頻度と比較するには、次を監査する。

\begin{enumerate}
\def\labelenumi{\arabic{enumi}.}
\tightlist
\item
  各開始点がほとんど確実に有限時間で次の開始面へ戻る。
\item
  1周期で結果レジスターが高々1回だけ確定する。
\item
  比較境界への複数回交差を重複計数しない。
\item
  基底、設定、結果に応じて失敗周期を捨てない。
\item
  開始数、設定生成数、局所記録数、比較完了数、外部記録数、再帰数を照合する。
\item
  Bell側で整合しない周期を観測後に捨てない。
\item
  外部記録と不要情報を含む拡大全系の写像を1対1に保つ。
\end{enumerate}

異なる結果履歴を外部記録ごと同じ点へ押しつぶす多対1写像は Hamiltonian ではない。結果情報は外部記録、不要情報モード、仕事源、環境のいずれかへ残す必要がある \cite{landauer1961,bennett1982}。

\hypertarget{ux9078ux629eux5668ux5c65ux6b74ux6bd4ux8f03ux5668ux53cdux5fa9ux5468ux671fux306e-hamiltonian-ux5019ux88dc}{%
\chapter{選択器、履歴比較器、反復周期の Hamiltonian 候補}\label{ux9078ux629eux5668ux5c65ux6b74ux6bd4ux8f03ux5668ux53cdux5fa9ux5468ux671fux306e-hamiltonian-ux5019ux88dc}}

\begin{statusbox}
位置づけ：有限基底選択器とBell履歴比較器に必要な正準基底混合、作用加算、有限幅比較、記録、逆計算の候補 Hamiltonian を示す。全構成要素を連結した自律 Hamiltonian、目的の不変母測度、Bell整合支持と共通基準密度は未導出である。
\end{statusbox}

\hypertarget{ux62e1ux5927ux5168ux7cfbux3068ux30a8ux30cdux30ebux30aeux30fcux53ceux652f}{%
\section{拡大全系とエネルギー収支}\label{ux62e1ux5927ux5168ux7cfbux3068ux30a8ux30cdux30ebux30aeux30fcux53ceux652f}}

1試行を担う有限部分を

\[
\begin{aligned}
H_{\rm fin}
={}&
H_{\rm particles}
+
H_{\rm ph}
+
H_{\rm prep}
+
H_{\rm program}
+
H_{\rm selector}
\\
&+
H_{\rm analyzer}
+
H_{\rm pointer}
+
H_{\rm return}
+
H_{\rm compare}
+
H_{\rm memory}
+
H_{\rm clk}
\end{aligned}
\]

とする。外部自由度と仕事源を含め、

\[
H_{\rm all}
=
H_{\rm fin}
+
H_{\rm ext}
+
\varepsilon_{\rm ext}H_{\rm link}
+
H_{\rm work}
\]

とする。\(H_{\rm all}\) は自律 Hamiltonian であり、有限部分だけのエネルギーは

\[
\dot E_{\rm fin}
=
J_{\rm in}
-
J_{\rm out}
+
P_{\rm ctrl}
\]

と変化し得る。測定窓では外部結合を小さくし、有限閉鎖系の正準写像として計算する。試行間では、外部記録、不要情報の移送、源と補助レジスターの再準備に弱開放流路を用い得る。

\hypertarget{ux6642ux8a08ux3068ux30d7ux30edux30b0ux30e9ux30e0ux7a93}{%
\section{時計とプログラム窓}\label{ux6642ux8a08ux3068ux30d7ux30edux30b0ux30e9ux30e0ux7a93}}

時計作用角を \((I_{\rm clk},\vartheta_{\rm clk})\) とし、

\[
H_{\rm clk}
=
\Omega_{\rm clk}I_{\rm clk}
\]

とする。各相互作用を、互いに分離した滑らかな周期窓 \(g_j(\vartheta_{\rm clk})\) で作動させる。

測定プログラム \(M\) は、正準レジスターの互いに分離した安定領域として実装する。相互作用を

\[
H_j
=
g_j(\vartheta_{\rm clk})
\chi_M(Q_M)
G_j
\]

とすれば、滑らかな窓 \(\chi_M\) により有限基底測定とBell測定を分岐できる。\(M\) を外部の非 Hamiltonian 命令として扱わず、開始面 \(\Sigma_0\) の変数へ含める。

\hypertarget{ux56faux5b9aux57faux5e95ux6df7ux5408ux306eux6b63ux6e96ux6027}{%
\section{固定基底混合の正準性}\label{ux56faux5b9aux57faux5e95ux6df7ux5408ux306eux6b63ux6e96ux6027}}

有限位相担体の実座標を \(Q,P\in\mathbb R^L\) とする。複素ユニタリ行列を \(U=A+iB\) と分けると、実表示

\[
\mathcal U_{\mathbb R}
=
\begin{pmatrix}
A&-B\\
B&A
\end{pmatrix}
\]

は直交かつシンプレクティックである。従って

\[
\begin{pmatrix}
Q'\\
P'
\end{pmatrix}
=
\mathcal U_{\mathbb R}
\begin{pmatrix}
Q\\
P
\end{pmatrix}
\]

は正準変換であり、複素表示では \(b'=Ub\) となる。

任意の有限 \(U\) は、2モード回転と位相回転の有限列へ分解できる。各回転を時計窓で作動させれば、有限正準回路として実装できる。ただし、単粒子側で同じ状態を異なる基底へ測るには、準備分布がプログラム \(U\) に依存しないことが別に必要である。

\hypertarget{ux9078ux629eux5668ux632fux52d5ux5b50}{%
\section{選択器振動子}\label{ux9078ux629eux5668ux632fux52d5ux5b50}}

選択器の正準対を作用角変数

\[
\left(
J_{\rm sel},
\vartheta_{\rm sel}
\right)
\]

とし、自由項を

\[
H_{\rm sel}
=
\omega_{\rm sel}J_{\rm sel}
\]

とする。固定作用セクター \(J_{\rm sel}=J_*\) では角が等速回転する。Haar一様角は自由回転で保存されるが、測定結合と再準備を含む全周期で条件付き一様性が保存されることは、この自由項だけから従わない。

角の切断を避けるには、円周上の基準点を測定窓の外へ置く。実装では \(\vartheta_{ m sel}/(2\pi)\) の周期的な滑らか近似 \(f_w(\vartheta_{ m sel})\) を用い、切断近傍の幅 \(w\) を有限比較誤差へ含める。

\hypertarget{ux4f5cux7528ux8aadux51faux3057ux3068ux7d2fux7a4dux30ecux30b8ux30b9ux30bfux30fc}{%
\section{作用読出しと累積レジスター}\label{ux4f5cux7528ux8aadux51faux3057ux3068ux7d2fux7a4dux30ecux30b8ux30b9ux30bfux30fc}}

測定基底混合後の局所作用を

\[
I_k
=
\frac12
\left[
\left(Q'_k\right)^2
+
\left(P'_k\right)^2
\right]
=
\mathcal J_0
\left|b'_k\right|^2
\]

とする。累積レジスターを \((Q_S,P_S)\) とし、\(k\) 番目の加算窓で

\[
H_{{\rm add},k}
=
g_k(\vartheta_{\rm clk})
P_S I_k
\]

を作動させる。パルス面積を1に取れば、自由発展を無視する理想極限で

\[
Q_S^{\rm out}
=
Q_S^{\rm in}
+
I_k,
\qquad
P_S^{\rm out}
=
P_S^{\rm in}
\]

となる。

位相担体への反作用は \(P_S\) に比例する。\(P_S=0\) の理想入口では読出し反作用が零であり、加算後も \(P_S=0\) が保たれる。他の Hamiltonian が加算窓中に \(P_S\) を生成しないことが必要である。

全作用を格納するレジスター \((Q_T,P_T)\) に対しても、

\[
H_{\rm total}
=
g_T(\vartheta_{\rm clk})
P_T
I_{\rm ph}
\]

を用いられる。選択閾値レジスター \((Q_U,P_U)\) へ

\[
H_u
=
g_u(\vartheta_{\rm clk})
P_U
I_{\rm ph}
f_w(\vartheta_{\rm sel})
\]

を作用させれば、理想的に \(Q_U=I_{\rm ph}\vartheta_{\rm sel}/(2\pi)\) を作る。角切断近傍と \(P_U\neq0\) の反作用を有限幅誤差へ含める。

\hypertarget{ux6709ux9650ux5e45ux6bd4ux8f03ux3068ux552fux4e00ux7d50ux679cux30e9ux30c3ux30c1}{%
\section{有限幅比較と唯一結果ラッチ}\label{ux6709ux9650ux5e45ux6bd4ux8f03ux3068ux552fux4e00ux7d50ux679cux30e9ux30c3ux30c1}}

累積値 \(Q_S\) と閾値 \(Q_U\) の差を、滑らかな単調関数 \(F_w\) で読む。判定レジスター \((Q_{D,k},P_{D,k})\) に

\[
H_{{\rm cmp},k}
=
g_{{\rm cmp},k}
(\vartheta_{\rm clk})
P_{D,k}
F_w
\left(
Q_S-Q_U
\right)
\]

を作用させると、\(Q_{D,k}\) は比較結果に応じて移動する。\(F_w\) は幅 \(w\) の遷移領域を持ち、\(w\to0\) で符号関数へ近づく。

最初の正の比較だけを結果にするには、空ラッチ領域の滑らかな窓 \(\chi_0(Q_L)\) と結果別ポインター生成子 \(G_{L,k}\) を用い、

\[
H_{{\rm latch},k}
=
g_{{\rm latch},k}
(\vartheta_{\rm clk})
\chi_0(Q_L)
\chi_+(Q_{D,k})
G_{L,k}
\]

とする。時計窓を \(k=1,\ldots,L\) の順に分離すれば、最初の正判定後は \(\chi_0(Q_L)\) が零となり、後の結果窓を抑制する候補になる。

この式は滑らかな一方向論理の候補であり、全ての初期条件で完全な一熱安定性、増幅、二重結果の不可能性を証明するものではない。有限幅での二重作動率と無結果率を検証する必要がある。

\hypertarget{ux7d50ux679cux8a18ux9332ux3068ux9006ux8a08ux7b97}{%
\section{結果記録と逆計算}\label{ux7d50ux679cux8a18ux9332ux3068ux9006ux8a08ux7b97}}

ラッチ結果を空の外部記録セルへ可逆にコピーする。記録コピー後、次を逆順に行う。

\begin{enumerate}
\def\labelenumi{\arabic{enumi}.}
\tightlist
\item
  ラッチ前の判定レジスターを逆計算する。
\item
  閾値と累積加算を逆計算する。
\item
  基底混合 \(U\) を逆実行する。
\item
  補助レジスターを基準領域へ戻す。
\end{enumerate}

加算の逆は \(-H_{{\rm add},k}\)、閾値計算の逆は \(-H_u\)、基底混合の逆は \(U^\dagger\) である。逆計算に必要な制御履歴を時計または不要情報レジスターへ保持する。

外部記録を含む全状態を同じ点へ戻すことはできない。異なる結果情報は外部記録、不要情報モード、仕事源、環境のどこかへ残る必要がある \cite{landauer1961,bennett1982}。

\hypertarget{bellux6e90ux3068ux5c40ux6240ux5206ux6790ux5668}{%
\section{Bell源と局所分析器}\label{bellux6e90ux3068ux5c40ux6240ux5206ux6790ux5668}}

Bell左右モードを \(z^A_{\mu r},z^B_{\nu r}\) とする。理想反対称源は

\[
\Xi_0
=
\sqrt{
\frac{\mathcal K}{2}
}
\left[
e_+e_-^{\mathsf T}
-
e_-e_+^{\mathsf T}
\right]
\]

である。2つの直交源チャンネルを

\[
z^{A,(1)}=a_0e_+,
\qquad
z^{B,(1)}=b_0e_-,
\]

\[
z^{A,(2)}=a_0e_-,
\qquad
z^{B,(2)}=-b_0e_+
\]

と準備すれば、係数の選択により \(\Xi_0\) を構成できる。

A側の局所回転生成子を

\[
G_A
=
\sum_r
\left(
Q^A_{+r}P^A_{-r}
-
Q^A_{-r}P^A_{+r}
\right)
\]

とする。\(H_A=\dot\alpha_xG_A\) はA側だけを \(R(\alpha_x)\) で回転する。B側も同様であり、

\[
\left\{
G_A,G_B
\right\}
=
0
\]

である。設定レジスター \(x,y\) はそれぞれの局所窓だけを制御する。

\hypertarget{ux623bux308aux4f1dux64adux3068ux679dux4f5cux7528ux8aadux51faux3057}{%
\section{戻り伝播と枝作用読出し}\label{ux623bux308aux4f1dux64adux3068ux679dux4f5cux7528ux8aadux51faux3057}}

局所出力から共通未来比較領域までの線形伝播を

\[
z_{{\rm ret},\mu r}^{X}(t)
=
\sum_{\nu,s}
\int
G_{\mu r,\nu s}^{X}(t-t')
z_{{\rm out},\nu s}^{X}(t')
\,dt'
\]

と書く。局所有限伝播には \(G^X(t)=0\)、\(t<\tau_X\) が必要である。

局所記録 \(A,B\) に対応する枝振幅 \(\Xi_{AB}^{xy}\) の実部と虚部を比較レジスター \((Q_{\rm R},P_{\rm R})\)、\((Q_{\rm I},P_{\rm I})\) へ写す候補は

\[
H_{\rm read}
=
g_{\rm read}(\vartheta_{\rm clk})
\left[
P_{\rm R}
\operatorname{Re}
\Xi_{AB}^{xy}
+
P_{\rm I}
\operatorname{Im}
\Xi_{AB}^{xy}
\right]
\]

である。\(P_{\rm R}=P_{\rm I}=0\) の理想入口では入力モードへの読出し反作用が零である。パルス面積を \(\Gamma\) とすると、

\[
Q_{\rm R}^{\rm out}
=
\Gamma
\operatorname{Re}
\Xi_{AB}^{xy},
\qquad
Q_{\rm I}^{\rm out}
=
\Gamma
\operatorname{Im}
\Xi_{AB}^{xy}
\]

となる。これらから

\[
K_{AB}^{xy}
=
\left|
\Xi_{AB}^{xy}
\right|^2
\]

と4つの累積区間境界を可逆算術レジスターへ計算する。比較器は局所結果の後に働き、結果を生成しない。

\hypertarget{bellux5883ux754cux89d2ux3068ux6574ux5408ux30ecux30b8ux30b9ux30bfux30fc}{%
\section{Bell境界角と整合レジスター}\label{bellux5883ux754cux89d2ux3068ux6574ux5408ux30ecux30b8ux30b9ux30bfux30fc}}

Bell選択器角を \(\vartheta_{\rm B}\) とし、

\[
u_{\rm B}
=
\frac{\vartheta_{\rm B}}{2\pi}
\mathcal K
\]

を閾値レジスターへ写す。局所記録済みの \((A,B)\) に対応する区間境界を \((T_{AB}^-,T_{AB}^+)\) とする。2つの有限幅比較器で

\[
T_{AB}^-
\leq
u_{\rm B}
<
T_{AB}^+
\]

を検査し、整合レジスターへ写す。

このレジスターは、既に記録された結果を変更しない。物理的完結履歴が整合領域だけに支持を持つことを表す終端自由度である。通常の前向き初期値問題で不整合レジスターを作った後にその周期を捨てるなら、事後選別になる。

\hypertarget{ux57faux6e96ux30bbux30afux30bfux30fcux5bc6ux5ea6ux306eux4f4dux7f6eux3065ux3051}{%
\section{基準セクター密度の位置づけ}\label{ux57faux6e96ux30bbux30afux30bfux30fcux5bc6ux5ea6ux306eux4f4dux7f6eux3065ux3051}}

Bell整合区間の長さは \(K_{AB}^{xy}\) である。しかし結果セクターの整合前基準密度 \(q_{AB}^{xy}\) は、上の算術回路から決まらない。\(q_{AB}^{xy}\) は次を含む。

\begin{enumerate}
\def\labelenumi{\arabic{enumi}.}
\tightlist
\item
  局所結果形成後のセクター密度。
\item
  時計面の正方向流束。
\item
  境界制約の余面積 Jacobian。
\item
  比較に直接参加しない付随自由度の体積。
\item
  完結履歴の解多重度。
\item
  結果窓と比較窓の有限幅。
\end{enumerate}

正規化された各結果セクターを正準写像で押し出しても、各セクターの総質量は保存される。正準角を追加するだけで \(q_{AB}^{xy}\) が共通になるわけではない。共通性は母測度のBell条件付き部分に必要な独立仮定である。

\hypertarget{ux5468ux671fux306eux5019ux88dcux9806ux5e8f}{%
\section{1周期の候補順序}\label{ux5468ux671fux306eux5019ux88dcux9806ux5e8f}}

統一周期の候補は次の順序である。

\begin{enumerate}
\def\labelenumi{\arabic{enumi}.}
\tightlist
\item
  時計が開始面 \(\Sigma_0\) を正方向に横切る。
\item
  源、プログラム、設定、選択器、空の補助レジスターを準備する。
\item
  位相担体を共通2次 Hamiltonian で伝播させる。
\item
  有限基底プログラムでは、\(U\)、累積加算、閾値比較、ラッチ、外部記録を実行する。
\item
  Bellプログラムでは、局所分析、局所結果記録、戻り伝播、枝作用計算、履歴整合比較を実行する。
\item
  結果を外部記録に残し、内部計算を可能な範囲で逆実行する。
\item
  消せない情報を不要情報または外部自由度へ移す。
\item
  時計が次の開始面へ戻る。
\end{enumerate}

この順序は周期写像 \(\mathcal R\) の設計図である。全項を1つの自律 Hamiltonian へ連結し、\(\mathcal R_*\mu_*=\mu_*\) とエルゴード性を示した完成模型ではない。

\hypertarget{ux6709ux9650ux5e45ux3068ux53cdux4f5cux7528}{%
\section{有限幅と反作用}\label{ux6709ux9650ux5e45ux3068ux53cdux4f5cux7528}}

理想パルス \(H_j\) と同時に自由項 \(H_0\) が働くとする。相互作用表示の先頭補正は概略

\[
\delta\mathcal U_j
=
O
\left(
\tau_j
\|H_0\|
\right)
+
O
\left(
\tau_j^2
\left\|
[H_j,H_0]
\right\|
\right)
\]

である。古典系では交換子を対応する Poisson作用素の交換子として読む。読出しレジスターの初期共役運動量が零でなければ、入力モードへの反作用が1次で生じる。

有限幅検証では次を別に測る。

\begin{enumerate}
\def\labelenumi{\arabic{enumi}.}
\tightlist
\item
  作用加算誤差。
\item
  選択角切断近傍の質量。
\item
  比較境界近傍の質量。
\item
  二重ラッチ率と無結果率。
\item
  結果別の記録失敗率。
\item
  Bell整合判定の誤り率。
\item
  逆計算後の補助レジスター残差。
\item
  次周期開始面への再帰誤差。
\end{enumerate}

\hypertarget{ux672aux5c0eux51faux4e8bux9805}{%
\section{未導出事項}\label{ux672aux5c0eux51faux4e8bux9805}}

本付録の候補 Hamiltonian からは、次は導かれない。

\begin{enumerate}
\def\labelenumi{\arabic{enumi}.}
\tightlist
\item
  目的の母測度 \(\mu_*\) の存在、一意性、不変性、エルゴード性。
\item
  選択器角の条件付き一様性を全周期で維持すること。
\item
  累積、比較、ラッチ、増幅が全初期条件で唯一結果を形成すること。
\item
  位相担体の選択結果が粒子の局所位置検出になること。
\item
  一般初期集団から反対称Bell源を準備すること。
\item
  Bell整合しない周期を事後選別なしで零測度にすること。
\item
  Bell基準セクター密度 \(q_{AB}^{xy}\) を4結果で共通にすること。
\item
  外部設定介入後にも同じBell統計を保つこと。
\item
  外部記録を保持しながら無期限に反復する完成周期。
\item
  平面内2出力を超える一般測定器と一般 Tsirelson原理。
\end{enumerate}

これらは補助模型内部の正準計算の厳密性とは別の、現行統一模型の準備、測定、境界、反復の課題である。

\begin{thebibliography}{99}

\addcontentsline{toc}{chapter}{参考文献}

\bibitem{bell1964} J. S. Bell, ``On the Einstein Podolsky Rosen Paradox,'' Physics Physique Fizika 1, 195--200 (1964). \url{https://doi.org/10.1103/PhysicsPhysiqueFizika.1.195}

\bibitem{chsh1969} J. F. Clauser, M. A. Horne, A. Shimony, and R. A. Holt, ``Proposed Experiment to Test Local Hidden-Variable Theories,'' Physical Review Letters 23, 880--884 (1969). \url{https://doi.org/10.1103/PhysRevLett.23.880}

\bibitem{nelson1966} E. Nelson, ``Derivation of the Schrödinger Equation from Newtonian Mechanics,'' Physical Review 150, 1079--1085 (1966). \url{https://doi.org/10.1103/PhysRev.150.1079}

\bibitem{guerra_morato1983} F. Guerra and L. M. Morato, ``Quantization of Dynamical Systems and Stochastic Control Theory,'' Physical Review D 27, 1774--1786 (1983). \url{https://doi.org/10.1103/PhysRevD.27.1774}

\bibitem{yasue1981} K. Yasue, ``Stochastic Calculus of Variations,'' Journal of Functional Analysis 41, 327--340 (1981). \url{https://doi.org/10.1016/0022-1236(81)90079-3}

\bibitem{zambrini1986} J.-C. Zambrini, ``Stochastic Mechanics According to E. Schrödinger,'' Physical Review A 33, 1532--1548 (1986). \url{https://doi.org/10.1103/PhysRevA.33.1532}

\bibitem{wharton2010} K. B. Wharton, ``Time-Symmetric Boundary Conditions and Quantum Foundations,'' Symmetry 2, 272--283 (2010). \url{https://doi.org/10.3390/sym2010272}

\bibitem{wharton_argaman2020} K. B. Wharton and N. Argaman, ``Colloquium: Bell's Theorem and Locally Mediated Reformulations of Quantum Mechanics,'' Reviews of Modern Physics 92, 021002 (2020). \url{https://doi.org/10.1103/RevModPhys.92.021002}

\bibitem{hall2010} M. J. W. Hall, ``Local Deterministic Model of Singlet State Correlations Based on Relaxing Measurement Independence,'' Physical Review Letters 105, 250404 (2010). \url{https://doi.org/10.1103/PhysRevLett.105.250404}

\bibitem{leifer_pusey2017} M. S. Leifer and M. F. Pusey, ``Is a Time Symmetric Interpretation of Quantum Theory Possible without Retrocausality?,'' Proceedings of the Royal Society A 473, 20160607 (2017). \url{https://doi.org/10.1098/rspa.2016.0607}

\bibitem{wood_spekkens2015} C. J. Wood and R. W. Spekkens, ``The Lesson of Causal Discovery Algorithms for Quantum Correlations,'' New Journal of Physics 17, 033002 (2015). \url{https://doi.org/10.1088/1367-2630/17/3/033002}

\bibitem{ford1965} G. W. Ford, M. Kac, and P. Mazur, ``Statistical Mechanics of Assemblies of Coupled Oscillators,'' Journal of Mathematical Physics 6, 504--515 (1965). \url{https://doi.org/10.1063/1.1704304}

\bibitem{mori1965} H. Mori, ``Transport, Collective Motion, and Brownian Motion,'' Progress of Theoretical Physics 33, 423--455 (1965). \url{https://doi.org/10.1143/PTP.33.423}

\bibitem{zwanzig1973} R. Zwanzig, ``Nonlinear Generalized Langevin Equations,'' Journal of Statistical Physics 9, 215--220 (1973). \url{https://doi.org/10.1007/BF01008729}

\bibitem{jamison1974} B. Jamison, ``Reciprocal Processes,'' Zeitschrift für Wahrscheinlichkeitstheorie und Verwandte Gebiete 30, 65--86 (1974). \url{https://doi.org/10.1007/BF00532864}

\bibitem{doob1957} J. L. Doob, ``Conditional Brownian Motion and the Boundary Limits of Harmonic Functions,'' Bulletin de la Société Mathématique de France 85, 431--458 (1957). \url{https://doi.org/10.24033/bsmf.1495}

\bibitem{landauer1961} R. Landauer, ``Irreversibility and Heat Generation in the Computing Process,'' IBM Journal of Research and Development 5, 183--191 (1961). \url{https://doi.org/10.1147/rd.53.0183}

\bibitem{bennett1982} C. H. Bennett, ``The Thermodynamics of Computation: A Review,'' International Journal of Theoretical Physics 21, 905--940 (1982). \url{https://doi.org/10.1007/BF02084158}

\bibitem{wallstrom1994} T. C. Wallstrom, ``Inequivalence between the Schrödinger Equation and the Madelung Hydrodynamic Equations,'' Physical Review A 49, 1613--1617 (1994). \url{https://doi.org/10.1103/PhysRevA.49.1613}

\bibitem{price_wharton2023} H. Price and K. Wharton, ``Bell Correlations as Selection Artefacts,'' arXiv:2309.10969v3 (2024). \url{https://arxiv.org/abs/2309.10969}

\bibitem{price_wharton2024} H. Price and K. Wharton, ``A Mechanism for Entanglement?,'' arXiv:2406.04571v1 (2024). \url{https://arxiv.org/abs/2406.04571}

\bibitem{argaman2010} N. Argaman, ``Bell's Theorem and the Causal Arrow of Time,'' American Journal of Physics 78, 1007--1013 (2010). \url{https://doi.org/10.1119/1.3456564}

\bibitem{hossenfelder_palmer2020} S. Hossenfelder and T. Palmer, ``Rethinking Superdeterminism,'' Frontiers in Physics 8, 139 (2020). \url{https://doi.org/10.3389/fphy.2020.00139}

\bibitem{thooft2016} G. 't Hooft, The Cellular Automaton Interpretation of Quantum Mechanics, Springer (2016). \url{https://doi.org/10.1007/978-3-319-41285-6}

\bibitem{leonard2014} C. Léonard, ``A Survey of the Schrödinger Problem and Some of Its Connections with Optimal Transport,'' Discrete and Continuous Dynamical Systems A 34, 1533--1574 (2014). \url{https://doi.org/10.3934/dcds.2014.34.1533}

\bibitem{chen_georgiou_pavon2016} Y. Chen, T. T. Georgiou, and M. Pavon, ``On the Relation between Optimal Transport and Schrödinger Bridges: A Stochastic Control Viewpoint,'' Journal of Optimization Theory and Applications 169, 671--691 (2016). \url{https://doi.org/10.1007/s10957-015-0803-z}

\bibitem{rauch_tung_striebel1965} H. E. Rauch, F. Tung, and C. T. Striebel, ``Maximum Likelihood Estimates of Linear Dynamic Systems,'' AIAA Journal 3, 1445--1450 (1965). \url{https://doi.org/10.2514/3.3166}

\bibitem{fuchs_goldt_seifert2016} J. Fuchs, S. Goldt, and U. Seifert, ``Stochastic Thermodynamics of Resetting,'' Europhysics Letters 113, 60009 (2016). \url{https://doi.org/10.1209/0295-5075/113/60009}

\bibitem{evans_majumdar_schehr2020} M. R. Evans, S. N. Majumdar, and G. Schehr, ``Stochastic Resetting and Applications,'' Journal of Physics A: Mathematical and Theoretical 53, 193001 (2020). \url{https://doi.org/10.1088/1751-8121/ab7cfe}

\bibitem{knorst_lopes2024} J. Knorst and A. O. Lopes, ``On the Quantum Guerra--Morato Action Functional,'' Journal of Mathematical Physics 65, 082102 (2024). \url{https://doi.org/10.1063/5.0207422}

\bibitem{wilson_et_al2021} J. T. Wilson, V. Borovitskiy, A. Terenin, P. Mostowsky, and M. P. Deisenroth, ``Pathwise Conditioning of Gaussian Processes,'' Journal of Machine Learning Research 22, 1--47 (2021). \url{https://jmlr.org/papers/v22/20-1260.html}

\bibitem{leonard_roelly_zambrini2014} C. Léonard, S. Rœlly, and J.-C. Zambrini, ``Reciprocal Processes. A Measure-Theoretical Point of View,'' Probability Surveys 11, 237--269 (2014). \url{https://doi.org/10.1214/13-PS220}

\bibitem{marchiori_deaguiar2011} M. A. Marchiori and M. A. M. de Aguiar, ``Energy Dissipation Via Coupling With a Finite Chaotic Environment,'' Physical Review E 83, 061112 (2011). \url{https://doi.org/10.1103/PhysRevE.83.061112}

\bibitem{heslot1985} A. Heslot, ``Quantum Mechanics as a Classical Theory,'' Physical Review D 31, 1341--1348 (1985). \url{https://doi.org/10.1103/PhysRevD.31.1341}

\bibitem{briggs_eisfeld2012} J. S. Briggs and A. Eisfeld, ``Coherent Quantum States from Classical Oscillator Amplitudes,'' Physical Review A 85, 052111 (2012). \url{https://doi.org/10.1103/PhysRevA.85.052111}

\bibitem{briggs_eisfeld2013} J. S. Briggs and A. Eisfeld, ``Quantum Dynamics Simulation with Classical Oscillators,'' Physical Review A 88, 062104 (2013). \url{https://doi.org/10.1103/PhysRevA.88.062104}

\bibitem{skinner2013} T. E. Skinner, ``Exact Mapping of the Quantum States in Arbitrary N-Level Systems to the Positions of Classical Coupled Oscillators,'' Physical Review A 88, 012110 (2013). \url{https://doi.org/10.1103/PhysRevA.88.012110}

\end{thebibliography}


\end{document}
